\documentclass[12pt,a4paper]{article}

% --- Πακέτα για Ελληνικά, Αγγλικά και Μαθηματικά ---
\usepackage[utf8]{inputenc}
\usepackage[english,greek]{babel}
\usepackage{amsmath}
\usepackage{amsfonts}
\usepackage{amssymb}
\usepackage{graphicx}
\usepackage{float} % Προαιρετικό, για να "κλειδώνεις" την εικόνα εκεί που θέλεις
\usepackage{subcaption}
\usepackage[a4paper, total={17cm, 25cm}]{geometry}
\usepackage{setspace}
\onehalfspacing  % Διάστιχο 1.5 (Προτεινόμενο για αναφορές)

\begin{document}

\begin{titlepage}
    \begin{center}
        \vspace*{1cm}
        
        % Λογότυπο (Αν έχεις το αρχείο ntua_logo.png, αλλιώς σχολίασε τη γραμμή)
        % \includegraphics[width=0.2\textwidth]{ntua_logo.png}\\[1cm]
        
        \Large
        \textbf{ΕΘΝΙΚΟ ΜΕΤΣΟΒΙΟ ΠΟΛΥΤΕΧΝΕΙΟ}\\
        \large
        ΣΧΟΛΗ ΗΛΕΚΤΡΟΛΟΓΩΝ ΜΗΧΑΝΙΚΩΝ ΚΑΙ ΜΗΧΑΝΙΚΩΝ ΥΠΟΛΟΓΙΣΤΩΝ\\[1.5cm]
        
        \vfill
        
        \Large
        \textbf{Μάθημα: Όραση Υπολογιστών}\\[0.5cm]
        
        \vfill
        
        \huge
        \textbf{Αναφορά: Εργαστήριο 1}\\[2cm]
        
        \vfill
        
        \Large
        \textbf{Στοιχεία Φοιτητών:}\\[0.5cm]
        \large
        Παναγόπουλος Μιχαήλ, Α.Μ: 03122061\\
        Φειδάκης Αθανάσιος, Α.Μ: 03122676\\[2cm]
        
        \vfill
        
        \large
        Αθήνα, Μάρτιος 2026
        
    \end{center}
\end{titlepage}

\newpage

\section*{Μέρος 1: Ανίχνευση Ακμών σε Γκρίζες Εικόνες}

\subsection*{1.1 Δημιουργία Εικόνων Εισόδου}

\subsubsection*{1.1.1 Εισαγωγή και Κανονικοποίηση Αρχικής Εικόνας}


\hspace{0.6cm}Η εικόνα "\foreignlanguage{english}{\textbf{edgetest\_26.png}}" φορτώνεται σε γκρι κλίμακα μέσω της εντολής \foreignlanguage{english}{\textbf{cv2.imread}} με το όρισμα \foreignlanguage{english}{\textbf{cv2.IMREAD\_GRAYSCALE}}, οι τιμές της αρχίκα κυμαίνονται στο διάστημα $[0, 255]$, με τύπο δεδομένων \foreignlanguage{english}{np.uint8} και κανονικοποιούνται (\foreignlanguage{english}{normalization}) στο διάστημα $[0, 1]$ με τύπο δεδομένων \foreignlanguage{english}{np.float64}.

Η κανονικοποίηση αυτή είναι απαραίτητη για την αριθμητική ευστάθεια, τη συνέπεια των ορίων κατά την εφαρμογή των μετέπειτα διαφορικών τελεστών και φίλτρων και την διευκόλυνση της προβολής της με χρήση της συνάρτησης \foreignlanguage{english}{\textbf{plt.imshow}}.

\begin{figure}[H] % Το [h] σημαίνει "here" (εδώ ακριβώς)
    \centering    % Στοίχιση στο κέντρο
    \includegraphics[scale=0.5]{figs_part1/fig1.png}
    \caption{Εικόνα \foreignlanguage{english}{Edgetest}} % Λεζάντα
    \label{fig:results} % Ετικέτα για αναφορά μέσα στο κείμενο (π.χ. Βλέπε Εικόνα \ref{fig:results})
\end{figure}

\subsubsection*{1.1.2 Μοντελοποίηση και Προσθήκη Θορύβου}
\hspace{0.6cm}Στην αρχική εικόνα $I_0$ προστίθεται λευκός \foreignlanguage{english}{Gaussian} θόρυβος. Επιλέγουμε δύο διαφορετικές τιμές για την τυπική απόκλιση του θορύβου ($\sigma_1$ = 0.1, $\sigma_2$ = 0.32), παράγοντας δύο θορυβώδη αντίγραφα της αρχικής εικόνας με \foreignlanguage{english}{PSNR} = 20 και \foreignlanguage{english}{PSNR} = 10 αντίστοιχα.
H τυπική απόκλιση του θορύβου $\sigma_n$, που μόλις παραθέσαμε, για μια επιθυμητή τιμή \foreignlanguage{english}{PSNR} υπολογίζεται ως:
\begin{equation}
\sigma_n = 10^{-\frac{\text{\foreignlanguage{english}{PSNR}}}{20}}
\end{equation}

Με την προσθήκη του θορύβου η εικόνα αλλοιώνεται με την εμφάνιση \foreignlanguage{english}{pixels} με τυχαία φωτεινότητα σε ορισμένες θέσεις της αρχικής εικόνας. Όσο πιο μεγάλη είναι η τιμή του σ τόσο πιο έντονη μεταβολή της φωτεινότητας μπορεί να προκύψει, καθώς ο θόρυβος λαμβάνει τιμές με μεγαλύτερη πιθανότητα απο τα άκρα της κατανομής. Τα αποτελέσματα της προσθήκης θορύβου είναι τα ακόλουθα:

\begin{figure}[H] % Το [h] σημαίνει "here" (εδώ ακριβώς)
    \centering    % Στοίχιση στο κέντρο
    \includegraphics[scale=0.6]{figs_part1/fig2.png}
    \caption{Εικόνα \foreignlanguage{english}{Edgetest} με Θόρυβο} % Λεζάντα
    \label{fig:results} % Ετικέτα για αναφορά μέσα στο κείμενο (π.χ. Βλέπε Εικόνα \ref{fig:results})
\end{figure}


\subsection*{1.2.1 Σχεδιασμός Φίλτρων και Διακριτοποίηση Κρουστικών Αποκρίσεων}

Για την ανίχνευση μεταβολών στην ένταση της εικόνας, χρησιμοποιούνται δύο θεμελιώδη φίλτρα:
\begin{enumerate}
    \item \textbf{Γκαουσιανό Φίλτρο ($G_\sigma$):} Λειτουργεί ως χαμηλοπερατό φίλτρο για την καταστολή του υψίσυχνου θορύβου (\foreignlanguage{english}{smoothing}). Ουσιαστικά απαλείφει τα \foreignlanguage{english}{peaks} του θορύβου ομογενοποιώντας τα με τα γειτονικά, μη θορυβώδη, \foreignlanguage{english}{pixels}. Έτσι αποφεύγονται τα ψευδή τοπικά μέγιστα του μέτρου του \foreignlanguage{english}{Gradient} της εικόνας. Η κρουστική απόκριση ορίζεται ως:
    \begin{equation}
    G_\sigma(x,y) = \frac{1}{2\pi\sigma^2} e^{-\frac{x^2+y^2}{2\sigma^2}}
    \end{equation}
    \item \textbf{Λαπλασιανή της Γκαουσιανής (\foreignlanguage{english}{LoG}):} Είναι ένας ισοτροπικός τελεστής που συνδυάζει την εξομάλυνση με τον υπολογισμό της δεύτερης παραγώγου. Ορίζεται ως:
    \begin{equation}
    \nabla^2 G_\sigma(x,y) = \left( \frac{x^2+y^2-2\sigma^2}{\sigma^4} \right) G_\sigma(x,y)
    \end{equation}
    Η χρησιμότητα του \foreignlanguage{english}{LoG} έγκειται στον εντοπισμό των μηδενισμών (\foreignlanguage{english}{zero-crossings}) της δεύτερης παραγώγου, που αντιστοιχούν σε τοπικά μέγιστα της κλίσης (ακμές).
\end{enumerate}

\subsubsection*{Αντιμετάθεση Παραγώγισης και Συνέλιξης}
Η Λαπλασιανή της εξομαλυμένης εικόνας δίνεται από:
\begin{equation}
L(x,y) = \nabla^2 [I(x,y) * G_\sigma(x,y)]
\end{equation}
Λόγω της ιδιότητας της συνέλιξης, ισχύει:
\begin{equation}
\nabla^2 [I * G_\sigma] = I * [\nabla^2 G_\sigma]
\end{equation}
\hspace{0.6cm}Η ιδιότητα αυτή επιτρέπει την αντιμετάθεση του τελεστή παραγώγισης από την θορυβώδη εικόνα στον πυρήνα του φίλτρου. Έτσι γίνεται χρήση ενός ενιαίου διαφορικού τελεστή ο οποίος ενσωματώνει τόσο την διαδικασία της εξομάλυνσης όσο και αυτήν της παραγώγισης. Με αυτόν τον τρόπο αποφεύγουμε τον υπολογισμό δύο συνελίξεων μειώνοντας το υπολογιστικό κόστος και τα σφάλματα προσέγγισης. Τα σφάλματα αυτά θα προέκυπταν σε πολύ μεγάλο βαθμό από την διακριτή προσέγγιση του  \foreignlanguage{english}{kernel} που υλοποιεί την παραγώγιση. Αντίθετα με το φίλτρο  \foreignlanguage{english}{LoG} αυτός ο ενιαίος  \foreignlanguage{english}{kernel} προκύπτει από την αναλυτική έκφραση (3).

\subsubsection*{Υλοποίηση σε \foreignlanguage{english}{Python}}
Για τη διακριτοποίηση των φίλτρων στην \foreignlanguage{english}{Python}, ακολουθήθηκε η εξής διαδικασία:
\begin{itemize}
    \item \textbf{Ορισμός μεγέθους πυρήνα:} Χρησιμοποιήθηκε η σχέση $n = 2 \cdot \lceil 3\sigma \rceil + 1$, ώστε ο πυρήνας να καλύπτει το ενεργό περιεχόμενο της κατανομής.
    \item \textbf{Κατασκευή πλέγματος:} Με τη συνάρτηση \texttt{\foreignlanguage{english}{np.meshgrid}}, δημιουργήθηκαν οι πίνακες συντεταγμένων $(x,y)$ στο διάστημα $[-(n-1)/2, (n-1)/2]$.
    \item \textbf{Υλοποίηση \foreignlanguage{english}{Gaussian}:} Χρησιμοποιήθηκε η \texttt{\foreignlanguage{english}{cv2.getGaussianKernel}}. Ο διδιάστατος πυρήνας προέκυψε από το εξωτερικό γινόμενο του μονοδιάστατου πίνακα που επέστρεψε η προηγούμενη συνάρτηση με τον ανάστροφό του. Ουσιαστικά εκμεταλλευόμαστε την ιδιότητα:
    \begin{equation}
    G_\sigma(x,y) = G_\sigma(x)G_\sigma(y)
    \end{equation}
    \item \textbf{Υλοποίηση \foreignlanguage{english}{LoG}:} Υλοποιήθηκε άμεσα μέσω της αναλυτικής έκφρασης (3) για μέγιστη ακρίβεια.
\end{itemize}

\subsection*{1.2.2 Προσέγγιση της Λαπλασιανής $L$ της Εξομαλυμένης Εικόνας}

\subsubsection*{Γραμμική Προσέγγιση ($L_1$): Συνέλιξη με τον Πυρήνα \foreignlanguage{english}{LoG}}
Εφαρμογή του γραμμικού χωρικά αμετάβλητου φίλτρου \foreignlanguage{english}{LoG} στην θορυβώδη εικόνα:
\begin{equation}
L_1(x,y) = \nabla^2(G_\sigma * I) = (\nabla^2 G_\sigma) * I = h(x,y) * I(x,y)
\end{equation}

Η υλοποίηση αυτή πραγματοποιείται ψηφιακά μέσω της συνέλιξης της αρχικής εικόνας με τον διακριτό πυρήνα που κατασκευάστηκε στο βήμα 1.2.1, χρησιμοποιώντας τη συνάρτηση \texttt{\foreignlanguage{english}{cv2.filter2D}}.

\subsubsection*{Μη-Γραμμική Προσέγγιση ($L_2$): Μορφολογική Λαπλασιανή}
Η δεύτερη προσέγγιση δίνεται από τον εξής τύπο

\begin{equation}
L_2(x,y) = (I_\sigma \oplus B)(x,y) + (I_\sigma \ominus B)(x,y) - 2I_\sigma(x,y)
\end{equation}

όπου $\oplus$ συμβολίζει τη μορφολογική διαστολή (\foreignlanguage{english}{dilation}) και $\ominus$ τη μορφολογική διάβρωση (\foreignlanguage{english}{erosion}).

\subsubsection*{Μαθηματική Θεμελίωση της Μορφολογικής Προσέγγισης}

\hspace{0.6cm}Η ποσότητα $(I_\sigma \oplus B) - I_\sigma$, γνωστή στην βιβλιογραφία ως \foreignlanguage{english}{edge}$^+$ αντιπροσωπεύει τo εξωτερικό περίγραμμα, ενώ η ποσότητα $I_\sigma - (I_\sigma \ominus B)$, γνωστή ως \foreignlanguage{english}{edge}$^-$, αντιπροσωπεύει το εσωτερικό περίγραμμα. Ωστόσο στην περίπτωση μας καλό είναι αυτές οι δύο ποσότητες να ερμηνευθούν ως  ````εξωτερική'''' και ````εσωτερική'''' κλίση αντίστοιχα. Η διαφορά αυτών των δύο κλίσεων, δηλαδή η διαφορά των δύο μεταβολών, μας δίνει μια εκτίμηση της δεύτερης παραγώγου:
\begin{equation}
L_2 = \underbrace{[(I_\sigma \oplus B) - I_\sigma]}_{\text{Άνω διαφορά}} - \underbrace{[I_\sigma - (I_\sigma \ominus B)]}_{\text{Κάτω διαφορά}} = (I_\sigma \oplus B) + (I_\sigma \ominus B) - 2I_\sigma
\end{equation}

Αυτή η τακτική προκύπτει ως λογική επέκταση της γνωστής προσέγγισης πεπερασμένων διαφορών $f(x+h) + f(x-h) - 2f(x)$ για την δεύτερη παράγωγο $f''(x)$ από την μια διάσταση.

\subsubsection*{Υλοποίηση σε \foreignlanguage{english}{Python}}
Για την υλοποίηση του $L_2$, χρησιμοποιήθηκαν οι συναρτήσεις \texttt{\foreignlanguage{english}{cv2.dilate}} και \texttt{\foreignlanguage{english}{cv2.erode}} με έναν πυρήνα \texttt{\foreignlanguage{english}{kernel}} που ορίζεται ως:
\\\\
\texttt{B =\foreignlanguage{english}{np.array}([[0, 1, 0],
        [1, 1, 1],
        [0, 1, 0]], \foreignlanguage{english}{dtype=np.uint8})}

Το αποτέλεσμα $L_2$ προκύπτει από την αριθμητική πρόσθεση και αφαίρεση των πινάκων που προκύπτουν.


\subsection*{1.2.3 Ανίχνευση Σημείων Μηδενισμού (\foreignlanguage{english}{Zero-Crossings})}

\subsubsection*{Θεωρητικό Υπόβαθρο}
\hspace{0.6cm}Οι ακμές μιας εικόνας αντιστοιχούν στα τοπικά μέγιστα της πρώτης παραγώγου της έντασης, τα οποία μαθηματικά ταυτίζονται με τους μηδενισμούς της δεύτερης παραγώγου. Στις δύο διαστάσεις, ο ισοτροπικός τελεστής που χρησιμοποιείται είναι η Λαπλασιανή $\nabla^2$. 

Η εύρεση των σημείων μηδενισμού της συνάρτησης $L(x,y)$ είναι προτιμότερη από την απλή κατωφλιωποίηση του μέτρου της κλίσης, διότι η εύρεση των σημείων μηδενισμού της δεύτερης παραγώγου υλοποιείται εύκολα και είναι πιο αποδοτική υπολογιστικά. Αντίθετα η εύρεση ακμών μέσω της ανίχνευσης τοπικών μεγίστων του \foreignlanguage{english}{Gradient} της εικόνας θα απαιτούσε πολλαπλές συγκρίσεις \foreignlanguage{english}{pixels} μεταξύ τους. Τέλος, η \foreignlanguage{english}{LoG} προσφέρει το πλεονέκτημα του σχηματισμού κλειστών και συνεχών καμπυλών.

\subsubsection*{Βήμα (α): Κατασκευή Δυαδικής Εικόνας Προσήμου}
\hspace{0.6cm}Αν βρισκόμασταν στον συνεχή χώρο τα σημεία των ακμών θα ήταν το σύνολο των σημείων $\{ (x,y) : L(x,y) = 0 \}$. Ωστόσο στο διακριτό χώρο, η ακμή βρίσκεται ανάμεσα στα \foreignlanguage{english}{pixels} όπου συμβαίνει η αλλαγή προσήμου. Για τον λόγω αυτό διαχωρίζουμε τα \foreignlanguage{english}{pixels} βάσει του προσήμου τους και σχηματίζουμε μια δυαδική εικόνα Χ.

\subsubsection*{Βήμα (β): Μορφολογικός Εντοπισμός Περιγράμματος}
\hspace{0.6cm}Για την εξαγωγή των σημείων μηδενισμού, χρησιμοποιούμε τον συμμετρικό μορφολογικό τελεστή περιγράμματος για δυαδικά σύνολα:
\begin{equation}
Y = (X \oplus B) \setminus (X \ominus B) = (X \oplus B) - (X \ominus B)
\end{equation}
\hspace{0.6cm}Συνεπώς, η διαφορά τους αφαιρεί το ````εσωτερικό'''' του συνόλου $X$ από το ````διευρυμένο'''' σύνολό του, αφήνοντας μόνο μια ζώνη πάχους ενός \foreignlanguage{english}{pixel} που διαζωρίζει τις τιμές $L \geq 0$ και $L < 0$. 

\subsubsection*{Υλοποίηση σε \foreignlanguage{english}{Python}}
\hspace{0.6cm}Στον κώδικα \foreignlanguage{english}{Python}, η αφαίρεση $(X \oplus B) - (X \ominus B)$ υλοποιείται με την αφαίρεση πινάκων. Το αποτέλεσμα είναι ένας δυαδικός χάρτης όπου οι τιμές 1 υποδεικνύουν την παρουσία ακμής.

\subsection*{1.2.4 Απόρριψη Ψευδών Ακμών σε Ομαλές Περιοχές}

\subsubsection*{Αναγκαιότητα του Επιπρόσθετου Ελέγχου}
\hspace{0.6cm}Παρόλο που ο εντοπισμός των σημείων μηδενισμού (\foreignlanguage{english}{zero-crossings}) της Λαπλασιανής παρέχει λεπτά και συνεχή περιγράμματα, ο τελεστής $\nabla^2$ εντοπίζει οποιαδήποτε αλλαγή προσήμου, ακόμη και αν αυτή προέρχεται από ανεπαίσθητες διακυμάνσεις της έντασης σε ομαλές (\foreignlanguage{english}{smooth}) περιοχές ή από υπολειπόμενο θόρυβο που δεν απορροφήθηκε από το Γκαουσιανό φίλτρο.

Για να διακρίνουμε τις ````πραγματικές'''' ακμές από τις ````ψευδείς'''', επιβάλλουμε ένα επιπλέον κριτήριο ισχύος. Μια έγκυρη ακμή δεν αρκεί μόνο να παρουσιάζει μηδενισμό της δεύτερης παραγώγου, αλλά πρέπει ταυτόχρονα να παρουσιάζει και σημαντική μεταβολή στην πρώτη παράγωγο.


\subsubsection*{Ο Ρόλος της Παραμέτρου $\theta_{edge}$}
Η παράμετρος $\theta_{edge} \in [0, 1]$ λειτουργεί ως κατώφλι της έντασης των ακμών:
\begin{enumerate}
    \item \textbf{Χαμηλές τιμές $\theta_{edge}$:} Επιτρέπουν τη διατήρηση ασθενών ακμών, αλλά αυξάνουν την πιθανότητα εμφάνισης θορύβου στην τελική εικόνα.
    \item \textbf{Υψηλές τιμές $\theta_{edge}$:} Καθαρίζουν την εικόνα από τον θόρυβο, αλλά ενδέχεται να προκαλέσου απώλεια τμημάτων των πραγματικών ακμών, ειδικά σε περιοχές με βαθμιαίες μεταβολές φωτεινότητας.
\end{enumerate}



\subsubsection*{Υλοποίηση σε \foreignlanguage{english}{Python}}
\hspace{0.6cm}Η υλοποίηση βασίστηκε στον υπολογισμό της κλίσης μέσω της \texttt{\foreignlanguage{english}{np.gradient}}, η οποία επιστρέφει τις μερικές παραγώγους κατά τον οριζόντιο και κατακόρυφο άξονα. Η τελική επιλογή των \foreignlanguage{english}{pixels} έγινε με τη χρήση του στοιχειακού λογικού τελεστή \foreignlanguage{english}{and} μεταξύ του δυαδικού πίνακα $Y$ και του πίνακα που προέκυψε από τα  \foreignlanguage{english}{pixels} που υπερέβησαν το κατώφλι.

\subsection*{1.3.1 Υπολογισμός Αληθινών Ακμών}

\hspace{0.6cm}Για την αξιολόγηση της απόδοσης των αλγορίθμων ανίχνευσης ακμών, είναι απαραίτητη η γνώση των «αληθινών» ακμών της αρχικής, καθαρής εικόνας $I_0$. Στην περίπτωση αυτή, οι ακμές $T$ προκύπτουν από την κατωφλίωποίηση του μορφολογικού περιγράμματος της εικόνας εισόδου.

\subsubsection*{Μαθηματική Επεξήγηση}
\hspace{0.6cm}Η πράξη $(I_0 \oplus B) - (I_0 \ominus B)$ είναι γνωστή ως Μορφολογική Κλίση και αντιστοιχίζεται στον απειροστικό λογισμό με την παράγωγο μιας συνάρτησης $f(x)$: $$f'(x) \approx \frac{f(x+h) - f(x-h)}{2h}$$
\hspace{0.6cm}Η διαφορά τους, μετράει τη μέγιστη διακύμανση της έντασης μέσα στο παράθυρο που ορίζει το δομικό στοιχείο $B$.Σε ομαλές περιοχές το μέγιστο και το ελάχιστο στη γειτονιά είναι σχεδόν ίδια. Άρα η διαφορά είναι κοντά στο $0$. (Μηδενική παράγωγος). Αντίθετα σε περιοχές με ακμή το μέγιστο θα είναι η τιμή της 'φωτεινής' πλευράς και το ελάχιστο η τιμή της 'σκοτεινής' πλευράς. Η διαφορά τους θα είναι μεγάλη, ίση με το άλμα της έντασης.
\\
\hspace{0.6cm}Σε αντίθεση με την κλασική προσέγγιση της διακριτής παραγώγου (π.χ. \foreignlanguage{english}{Sobel}) που υπολογίζει την κλίση σε μια συγκεκριμένη κατεύθυνση ($x$ ή $y$), η μορφολογική κλίση είναι ισοτροπική. Επειδή παίρνει το $max$ και το $min$ σε όλη τη γειτονιά $B$, ανιχνεύει την απότομη αλλαγή προς οποιαδήποτε κατεύθυνση ταυτόχρονα.

Αξίζει να σημειωθεί ότι η επιλογή του $\theta_{realedge}$ έγινε με τέτοιο τρόπο ώστε να απομονώνονται μόνο οι σημαντικές μεταβολές της $I_0$, δημιουργώντας έναν αξιόπιστο χάρτη αναφοράς για τον υπολογισμό των σφαλμάτων των υπόλοιπων μεθόδων.


\subsubsection*{Γιατί δεν απαιτείται Γκαουσιανή εξομάλυνση}
\hspace{0.6cm}Στην επεξεργασία της αρχικής εικόνας $I_0$, η παράλειψη του σταδίου της Γκαουσιανής εξομάλυνσης είναι όχι μόνο επιτρεπτή, αλλά και επιβεβλημένη γιατι ο θόρυβος πλέον απουσιάζει και η εφαρμογή ενός χαμηλοπερατου φίλτρου θα οδηγούσε σε χωρικό άπλωμα των ακμών και μείωση της οξύτητάς τους.

\subsection*{1.3.2 Ποσοτική Αξιολόγηση και Μετρικές Απόδοσης}

Για την αντικειμενική αξιολόγηση του αλγορίθμου, συγκρίνουμε το σύνολο των ανιχνευθέντων ακμών με το σύνολο των αληθινών ακμών μέσω των εξής μετρικών:

\begin{itemize}
    \item \textbf{(\foreignlanguage{english}{Recall} - $Pr(T|D)$):} Μετράει το ποσοστό των αληθινών ακμών που καταφέραμε να ανιχνεύσουμε. Υψηλό \foreignlanguage{english}{Recall} σημαίνει ότι δεν χάσαμε σημαντικές ακμές.
    \item \textbf{(\foreignlanguage{english}{Precision} - $Pr(D|T)$):} Μετράει πόσες από τις ακμές που ανιχνεύσαμε είναι όντως αληθινές. Υψηλό \foreignlanguage{english}{Precision} σημαίνει ότι ο αλγόριθμος δεν παρήγαγε ψευδείς ακμές (\foreignlanguage{english}{false positives}) λόγω θορύβου.
\end{itemize}

\subsubsection*{Σύνθετο Κριτήριο $C$ και Ισορροπία Μετρικών}
\hspace{0.6cm}Το κριτήριο ποιότητας $C = \frac{Pr(D|T) + Pr(T|D)}{2}$ αποτελεί τον αριθμητικό μέσο των δύο παραπάνω μετρικών. Η χρήση του είναι επιβεβλημένη διότι στην ανίχνευση ακμών υπάρχει πάντα ένα \foreignlanguage{english}{trade off} ανάμεσα στις δύο προαναφερθείσες μετρικές. Προσπαθώντας να αυξήσουμε την τιμή της μίας αναποφευκτα μειώνεται η τιμή της άλλης. Συνεπώς η βέλτιστη απόδοση του συστήματος επιτυγχάνεται όταν το κριτήριο $C$ μεγιστοποιείται, υποδεικνύοντας τις παραμέτρους που επιτυγχάνουν την καλύτερη δυνατή ισορροπία μεταξύ της πληρότητας της ανίχνευσης και της καθαρότητας από θόρυβο.

\subsection*{1.3.3 Μεταβολή Παραμέτρων και Χρήσιμα Συμπεράσματα}

1. \textbf{Σύγκριση Μεθόδων: Γραμμική \foreignlanguage{english}{vs} Μη-Γραμμική }

Παρατηρούμε ότι σχεδόν σε όλους τους συνδυασμούς $(\sigma, \theta)$, η Μη-Γραμμική μέθοδος υπερτερεί σε \foreignlanguage{english}{Precision}, \foreignlanguage{english}{Recall} και τελικό \foreignlanguage{english}{C-Score} (για τους πίνακες βλ. Παράρτημα). Από τα Σχήματα 3 και 4 βλέπουμε τις ακμές των μορφολογικών φίλτρων καλύτερα εντοπισμένες στον χώρο (μικρές ταλαντώσεις στα σχήματα Τ και Ο). 
\\\\
\textbf{Αιτιολόγηση}: 
Η γραμμική προσέγγιση  βασίζεται στη συνέλιξη με τον \foreignlanguage{english}{LoG} πυρήνα $h$ ($L_1 = h * I$). Η συνέλιξη, ως γραμμικός τελεστής, τείνει να διαχέει τον θόρυβο στις γειτονικές περιοχές. Αντίθετα $L_2$ χρησιμοποιεί μορφολογικά φίλτρα : $L_2 = I_{\sigma} \oplus B + I_{\sigma} \ominus B - 2I_{\sigma}$. Οι μη-γραμμικοί αυτοί τελεστές είναι πιο ανθεκτικοί στον θόρυβο και παρουσιάζουν καλύτερη απόκριση στις απότομες μεταβολές χωρίς να θολώνουν υπερβολικά τη δομή της εικόνας, οδηγώντας σε ακριβέστερο εντοπισμό των \foreignlanguage{english}{zero-crossings}.
\\\\
2.\textbf{Επίδραση της Παραμέτρου $\sigma$ (Χωρική Κλίμακα)} 

Η παράμετρος $\sigma$ ελέγχει το μέγεθος του \foreignlanguage{english}{Gaussian} πυρήνα ομαλοποίησης ($G_{\sigma}$). Μικρό $\sigma$ (π.χ. 1.0): Προσφέρει υψηλή ακρίβεια στη θέση της ακμής, αλλά δεν καταστέλλει επαρκώς τον θόρυβο. Αυτό οδηγεί σε χαμηλό \foreignlanguage{english}{Precision}, καθώς πολλές διακυμάνσεις του θορύβου εκλαμβάνονται λανθασμένα ως ακμές (\foreignlanguage{english}{false positives}). Μεγάλο $\sigma$ (π.χ. 3.0+) καταστέλλει αποτελεσματικά τον θόρυβο, αλλά προκαλεί έντονη θόλωση. Αυτό σβήνει τις λεπτές ακμές και μετατοπίζει τις υπάρχουσες, οδηγώντας σε πτώση του \foreignlanguage{english}{Recall} (πολλές αληθινές λεπτές ακμές δεν ανιχνεύονται). 

Στο σχήμα 3 και στην εικόνα με \foreignlanguage{english}{PSNR}=10 (σ=3) βλέπουμε πολύ πιο έντονη καταστολή των λευκών κουκίδων του θορύβου συγκριτικά με την αντίστοιχη εικόνα στο Σχήμα 4 (σ=2). Αντίθετα το χαμηλότερο σ οδηγεί \foreignlanguage{english}{blur} στις ακμές των σχημάτων Τ και Ο με αποτέλεσμα στο Σχήμα 4 να μην εμφανίζοντα οι ακανόνιστες δομές του Σχήματος 3. 

Για 20 \foreignlanguage{english}{dB} (χαμηλός θόρυβος), το βέλτιστο είναι 1.5, καθώς αρκεί μια ήπια εξομάλυνση. Για 10 \foreignlanguage{english}{dB} (υψηλός θόρυβος), η βέλτιστη τιμή ανεβαίνει στο 2.0, καθώς απαιτείται ισχυρότερο φιλτράρισμα για να επιζήσει η πληροφορία της ακμής.
\\
\\
3. \textbf{Επίδραση της Παραμέτρου $\theta_{edge}$}

Το $\theta_{edge}$ χρησιμοποιείται για την απόρριψη των \foreignlanguage{english}{zero-crossings} σε ομαλές περιοχές. Καθώς το $\theta$ αυξάνεται, το \foreignlanguage{english}{Precision} αυξάνεται σταθερά. Αυτό συμβαίνει γιατί το φίλτρο γίνεται πιο επιλεκτικό και κρατάει μόνο τις πολύ έντονες μεταβολές, απορρίπτοντας τις ````ψεύτικες'''' ακμές από τον θόρυβο. Ταυτόχρονα όμως, το \foreignlanguage{english}{Recall} μειώνεται, διότι ένα πολύ υψηλό κατώφλι απορρίπτει και αληθινές ακμές χαμηλής αντίθεσης. 

Βάσει των πειραμάτων, στο δεύτερο ζεύγος του Σχήματος 4 επιλέγεται μεγαλύτερο θ ώστε να καταστείλει πιο δραστικά τα ψευδή λευκά  \foreignlanguage{english}{dots} στο  \foreignlanguage{english}{backround} της εικόνας που προέρχονται από την μείωση του σ (συγκριτκικά με το Σχήμα 3). Αυτό ωστόσο ο οδηγεί σε περισσότερες ασυνέχειες στην περιφέρεια του κύκλου και του Τ. 

Στα 20 \foreignlanguage{english}{dB} το βέλτιστο είναι 0.20, ενώ στα 10 \foreignlanguage{english}{dB} ανεβαίνει στο 0.25. Αυτή η μετατόπιση είναι αναγκαία για να αντισταθμιστεί η μεγαλύτερη ένταση των τυχαίων διακυμάνσεων (θορύβου) στο σήμα.
\\
\\
4.\textbf{ Ανάλυση Βέλτιστων Αποτελεσμάτων (\foreignlanguage{english}{C-Score)}}
\\
Το κριτήριο $C$ είναι ο μέσος όρος \foreignlanguage{english}{Precision} και \foreignlanguage{english}{Recall} όπως αμαφέρθηκε ήδη.
\\
\foreignlanguage{english}{PSNR} = 20 \foreignlanguage{english}{dB} ($C_{max} = 0.932$): Επιτυγχάνεται με $\sigma=1.5, \theta=0.2, Non-Linear$. Εδώ η εικόνα έχει σχετικά καλό λόγο σήματος προς θόρυβο, επιτρέποντας μια μέση τιμή εξομάλυνσης και κατωφλίου που ισορροπεί το \foreignlanguage{english}{Recall} με το \foreignlanguage{english}{Precision}.
\\
\foreignlanguage{english}{PSNR} = 10 \foreignlanguage{english}{dB} ($C_{max} = 0.757$): Επιτυγχάνεται με $\sigma=2.0, \theta=0.25, Non-Linear$. Η πτώση του \foreignlanguage{english}{C-Score} (από 0.932 σε 0.757) οφείλεται στο ότι ο θόρυβος είναι πλέον τόσο έντονος που επικαλύπτει μέρος της δομικής πληροφορίας. Η βέλτιστη λύση μετατοπίζεται σε μεγαλύτερο $\sigma$ (για περισσότερο \foreignlanguage{english}{smoothing}) και υψηλότερο $\theta$ (για αυστηρότερο φιλτράρισμα), θυσιάζοντας λίγο \foreignlanguage{english}{Recall} για να διατηρηθεί ένα αξιοπρεπές \foreignlanguage{english}{Precision}.

\begin{figure}[H] % Το [h] σημαίνει "here" (εδώ ακριβώς)
    \centering    % Στοίχιση στο κέντρο
    \includegraphics[scale=0.5]{figs_part1/fig3.png}
    \caption{Ακμές \foreignlanguage{english}{Edgetest} με Ενδεικτικές Παρεμέτρους} % Λεζάντα
    \label{fig:results} % Ετικέτα για αναφορά μέσα στο κείμενο (π.χ. Βλέπε Εικόνα \ref{fig:results})
\end{figure}

\begin{figure}[H] % Το [h] σημαίνει "here" (εδώ ακριβώς)
    \centering    % Στοίχιση στο κέντρο
    \includegraphics[scale=0.5]{figs_part1/fig4.png}
    \caption{Ακμές \foreignlanguage{english}{Edgetest} με Βέλτιστες Παρεμέτρους} % Λεζάντα
    \label{fig:results} % Ετικέτα για αναφορά μέσα στο κείμενο (π.χ. Βλέπε Εικόνα \ref{fig:results})
\end{figure}

\subsection*{1.4.1 Εφαρμογή σε Πραγματική Εικόνα (\foreignlanguage{english}{ermoupoli.jpg})}

\hspace{0.6cm}Στο τελευταίο στάδιο του πρώτου μέρους, ελέγχεται η ικανότητα γενίκευσης του αλγορίθμου ανίχνευσης ακμών σε μία πραγματική εικόνα υψηλής πολυπλοκότητας. Η εικόνα \texttt{\foreignlanguage{english}{ermoupoli.jpg}} διακρίνεται από έντονες γεωμετρικές δομές (κτίρια, παράθυρα) και εμφανίζει λεπτά χαρακτηριστικά. Για την ανίχνευση ακμών επιλέχθηκαν τα μη γραμμικά φίλτρα καθώς στα προηγούμενα ερωτήματα διαπιστώσαμε μέσω της μετρικής \foreignlanguage{english}{C-score} ότι υπερέχει. Για μια πρώτη εφαρμογή επιλέχθηκαν οι τιμές $\sigma = 1.5$ και $\theta_{edge} = 0.2$ για τις παραμέτρους υποθέτοντας υπαρκτό αλλά μη αισθητό εγγενή θόρυβο.

\begin{figure}[H] % Το [h] σημαίνει "here" (εδώ ακριβώς)
    \centering    % Στοίχιση στο κέντρο
    \includegraphics[scale=0.7]{fig5.png}
    \caption{Ανίχνευση Ακμών στην εικόνα \foreignlanguage{english}{ermoupoli.jpg}} % Λεζάντα
    \label{fig:results} % Ετικέτα για αναφορά μέσα στο κείμενο (π.χ. Βλέπε Εικόνα \ref{fig:results})
\end{figure}

\subsection*{1.4.2 Διερεύνηση Παραμέτρων και Ποιοτική Αξιολόγηση}

\hspace{0.6cm}Στο σημείο αυτό εξετάζεται η ευαισθησία του ανιχνευτή ακμών ως προς τη χωρική κλίμακα $\sigma$ του \foreignlanguage{english}{Gaussian} φίλτρου και το κατώφλι έντασης $\theta_{edge}$ πάνω στην πραγματική εικόνα. Η κατανόηση του τρόπου με τον οποίο επενεργούν πάνω στην εικόνα θα επιτρέψουν την εκλογή του βέλτιστου (ποιοτικά) ζεύγους.

\subsection*{Επίδραση της Χωρικής Κλίμακας $\sigma$ (Σταθερό $\theta=0.2$)}
\hspace{0.6cm}Η παράμετρος $\sigma$ καθορίζει το βαθμό εξομάλυνσης της εικόνας πριν την παραγώγιση. Παρατηρώντας την πρώτη γραμμή του Σχήματος 6 μπορούμε να οδηγηθούμε στα εξής συμπεράσματα:

\textbf{Μικρό $\sigma$ (0.5):} Ο ανιχνευτής είναι εξαιρετικά ευαίσθητος στις λεπτομέρειες, αλλά καταγράφει πολύ θόρυβο και υφές (π.χ. λεπτομέρειες στα σκαλιά του κτηρίου, στην 'πριωνοτή' υφή του κορμού του φοίνικα και του φυλλώματος των δένδρων, στις εξοχές των παραθύρων) που μπορεί να μην αποτελούν κύριες ακμές. 

\textbf{Μεγάλο $\sigma$ (2.5 - 5.0):} Το φίλτρο λειτουργεί ως ισχυρό χαμηλοπερατό φίλτρο. Οι δευτερεύουσες ακμές εξαφανίζονται και παραμένουν μόνο οι κύριοι όγκοι των κτιρίων, τα βασικά περιγράμματα που λειτουργούν ως διαχωριστικά ανάμεσα στα αντικείμενα. Οι ακμές χάνουν την ακρίβεια εντοπισμού τους στον χώρο και οι γωνίες στρογγυλεύουν.


\subsection*{Επίδραση του Κατωφλίου $\theta_{edge}$ (Σταθερό $\sigma=1.5$)}
Το κατώφλι $\theta_{edge}$ καθορίζει ποιες από τις ανιχνευθείσες αλλαγές προσήμου (μηδενισμοί) θα θεωρηθούν έγκυρες ακμές βάσει του μέτρου της κλίσης τους. Παρατηρώντας την δεύτερη γραμμή του Σχήματος 6 μπορούμε να οδηγηθούμε στα εξής συμπεράσματα:

\textbf{Χαμηλό $\theta$ (0.1):} Ο χάρτης ακμών είναι πολύ πυκνός. Συμπεριλαμβάνονται ασθενείς μεταβολές, με αποτέλεσμα η εικόνα να εμφανίζει και τις πιο μικρές λεπτομέρειες, μικρές μεταβολές φωτεινότητας, συμπεριλαμβάνοντας αναπόφευκτα και ````οπτικό θόρυβο''''.

\textbf{Υψηλό $\theta$ (0.8):} Ο ανιχνευτής γίνεται πολύ επιλεκτικός. Μόνο οι ακμές με την υψηλότερη αντίθεση επιβιώνουν, οδηγώντας σε έναν φτωχή απεικόνιση με πολλές ασυνέχεις στις ανιχνευθείσες ακμές.

\begin{figure}[H] % Το [h] σημαίνει "here" (εδώ ακριβώς)
    %\centering    % Στοίχιση στο κέντρο
    \includegraphics[scale=0.4]{figs_part1/fig6.png}
    \caption{Μεταβολή παραμέτρων σ και θ για την εικόνα \foreignlanguage{english}{ermoupoli.jpg}} % Λεζάντα
    \label{fig:results} % Ετικέτα για αναφορά μέσα στο κείμενο (π.χ. Βλέπε Εικόνα \ref{fig:results})
\end{figure}


\subsection*{Βέλτιστο Αποτέλεσμα}

\hspace{0.6cm}Ως βέλτιστος συνδυασμός για την εικόνα προκρίνεται το ζεύγος \textbf{$\sigma=1.2$} και \textbf{$\theta_{edge}=0.2$}. Ο συνδυασμός αυτός επιτυγχάνει την καλύτερη ισορροπία καθώς η εξομάλυνση είναι επαρκής για να καθαρίσει την εικόνα χωρίς να αλλοιώσει τη γεωμετρία, ενώ το κατώφλι διατηρεί τη συνεκτικότητα των περιγραμμάτων χωρίς να επιτρέπει την εμφάνιση ψευδών ακμών.

\begin{figure}[H] % Το [h] σημαίνει "here" (εδώ ακριβώς)
    \centering    % Στοίχιση στο κέντρο
    \includegraphics[scale=0.8]{figs_part1/fig7.png}
    \caption{Βέλτιστο αποτέλεσμα ανίχνευσης στην εικόνα \foreignlanguage{english}{ermoupoli.jpg}} % Λεζάντα
    \label{fig:results} % Ετικέτα για αναφορά μέσα στο κείμενο (π.χ. Βλέπε Εικόνα \ref{fig:results})
\end{figure}

\subsection*{Σχόλιο}
\hspace{0.6cm}Οι παρατηρήσεις που ανέκυψαν φαίνεται να εναρμονίζονται πλήρως με το θεωρητικό διάγραμμα του Σχήματος 8. Το διάγραμμα αυτό πάρθηκε από το Κεφάλαιο 11 του βιβλίου ````\foreignlanguage{english}{P. Maragos, Images Analysis and Computer Vision}''''. Στο διάγραμμα φαίνεται πως μεγάλα σ επιτρέπουν την ανίχνευση μόνο παχιών ακμών (αντικείμενα μεγάλης κλίμακας) ολισθημένες ως προς την αρχική τους θέση, ενώ μικρά σ εντοπίζουν και πιο λεπτές ακμές (αντικείμενα μικρής κλίμακας) διατηρώντας την αρχική τους θέση στο χώρο. 

\begin{figure}[H] % Το [h] σημαίνει "here" (εδώ ακριβώς)
    \centering    % Στοίχιση στο κέντρο
    \includegraphics[scale=0.5]{figs_part1/figur7_5.png}
    \caption{Παράδειγμα ανιχνευθέντων ακμών για συνεχείς τιμές σ} % Λεζάντα
    \label{fig:results} % Ετικέτα για αναφορά μέσα στο κείμενο (π.χ. Βλέπε Εικόνα \ref{fig:results})
\end{figure}


\section*{Μέρος 2: Ανίχνευση Σημείων Ενδιαφέροντος (\foreignlanguage{english}{Interest Point Detection})}

\subsection*{2.1 Ανίχνευση Γωνιών}
\subsection*{\textbf{Θεωρητική Επεξήγηση: Μέθοδος \foreignlanguage{english}{Harris-Stephens}}}
\hspace{0.6cm}Ο αλγόριθμος \foreignlanguage{english}{Harris-Stephens} προσεγγίζει το πρόβλημα αξιολογώντας τη μεταβολή της έντασης κατά τη μετατόπιση ενός χωρικού παραθύρου. Στο πρώτο βήμα του αλγορίθμου ζητείται ο υπολογισμός των στοιχείων του δομικού τανυστή \foreignlanguage{english}{J}. Η μορφή του τανυστή αυτού προκύπτει απο την συνάρτηση διαφοράς αυτοσυσχέτισης για μια απειροστή μετατόπιση $(u, v)$ και για μια συνάρτηση βάρους $w(x,y)$ εκφράζεται ως: $$E(u, v) = \sum_{x,y} w(x,y) \left[ I_\sigma(x+u, y+v) - I_\sigma(x, y) \right]^2$$ 

Στην παραπάνω συνάρτηση εμπεριέχεται η εξομαλυμένη εικόνα $I_\sigma = G_\sigma * I$, αντί της αρχικής, για τη διασφάλιση συνεχούς διαφορισιμότητας και αποθορυβοποίησης, ενώ το εσωτερικό άθροισμα με το παράθυρο $w(x,y)$ αντιστοιχίζεται - στα επόμενα βήματα - στον συνεχή χώρο με συνέλιξη με τον πυρήνα ολοκλήρωσης $G_\rho(x,y)$. Χρησιμοποιώντας το ανάπτυγμα \foreignlanguage{english}{Taylor} πρώτης τάξης γύρω από το σημείο $(x, y)$, η μετατοπισμένη εικόνα προσεγγίζεται ρητά από τη γραμμική συνιστώσα:$$I_\sigma(x+u, y+v) \approx I_\sigma(x, y) + u \frac{\partial I_\sigma}{\partial x} + v \frac{\partial I_\sigma}{\partial y} = I_\sigma(x,y) + \begin{bmatrix} u & v \end{bmatrix} \begin{bmatrix} I_x \\ I_y \end{bmatrix}$$ 

Αντικαθιστώντας τη γραμμική αυτή προσέγγιση στη συνάρτηση σφάλματος $E(u,v)$ και εφαρμόζοντας στοιχειώδεις πράξεις γραμμικής άλγεβρας, η τετραγωνική μορφή γράφεται ως :$$E(u, v) \approx \sum_{x,y} w(x,y) \left( \begin{bmatrix} u & v \end{bmatrix} \begin{bmatrix} I_x \\ I_y \end{bmatrix} \right)^2 = \begin{bmatrix} u & v \end{bmatrix} \left( \sum_{x,y} w(x,y) \begin{bmatrix} I_x^2 & I_x I_y \\ I_x I_y & I_y^2 \end{bmatrix} \right) \begin{bmatrix} u \\ v \end{bmatrix}$$ 

Απομονώνοντας το εσωτερικό τετραγωνικό πίνακα και αντικαθιστώντας το άθροισμα με συνέλιξη βρίσκουμε τον τύπο του δομικού τανυστή $J(x, y) \in \mathbb{R}^{2 \times 2}$ 
$$J(x, y) = G_\rho * (\nabla I_\sigma \nabla I_\sigma^T) = \begin{bmatrix} G_\rho * (\frac{\partial I_\sigma}{\partial x})^2 & G_\rho * (\frac{\partial I_\sigma}{\partial x}\frac{\partial I_\sigma}{\partial y}) \\ G_\rho * (\frac{\partial I_\sigma}{\partial x}\frac{\partial I_\sigma}{\partial y}) & G_\rho * (\frac{\partial I_\sigma}{\partial y})^2 \end{bmatrix} = \begin{bmatrix} J_{1}(x,y) & J_{2}(x,y) \\ J_{2}(x,y) & J_{3}(x,y) \end{bmatrix}$$

Ως δεύτερο βήμα του αλγορίθμου ορίζεται η εύρεση των ιδιοδιανσυμάτων και ιδιοτιμών του τανυστή \foreignlanguage{english}{J}. Το μητρώο $J$ είναι πραγματικό, συμμετρικό και θετικά ημιορισμένο, γεγονός που εγγυάται ότι οι ιδιοτιμές του είναι πραγματικοί, μη αρνητικοί αριθμοί $\lambda_{\pm} \ge 0$. Τα ιδιοδιανύσματα $w_+, w_-$ του τανυστή αντιστοιχούν στις διευθλυνσεις μέγιστης και ελάχιστης μεταβολής της έντασης της εικόνας και οι αντίστοιχες ιδιοτιμές $\lambda_+, \lambda_-$ παρέχουν μια μέτρηση των μεταβολών αυτών. Οι τιμές των ιδιοτιμών κατηγοριοποιεί τα σημεία της εικόνας ως εξής:
\\\\
Ομαλή περιοχή, χωρίς σημαντική μεταβολή έντασης: $$\lambda_+ \approx 0, \ \lambda_- \approx 0$$
\\
Σημείο ακμής, όπου η μεταβολή συγκεντρώνεται σε μία κυρίαρχη διεύθυνση (το ιδιοδιάνυσμα που αντιστοιχεί στο $\lambda_+$):$$\lambda_+ \gg 0, \ \lambda_- \approx 0$$
\\
Γωνία ή κορυφή, το σήμα μεταβάλλεται έντονα προς όλες τις διευθύνσεις:$$\lambda_+ \gg 0, \ \lambda_- \gg 0$$

Το τελευταίο βήμα του αλγορίθμου αποτελεί η διατύπωση του κριτηρίου γωνιότητας που συμπυκνώνει υπολογιστικά αποδοτικά τις παραπάνω ιδιότητες των ιδιοτιμών. Με την εφαρμογή δύο συνθηκών (τοπικό μέγιστο και κατωφλιοποίηση) εξάγονται οι γωνίες.


\subsection*{Υλοποίηση σε \foreignlanguage{english}{Pyhton}}
\hspace{0.6cm}Η συνάρτηση \foreignlanguage{english}{CornersDetect} λαμβάνει ως όρισμα το \foreignlanguage{english}{plots} (με \foreignlanguage{english}{default} τιμή \foreignlanguage{english}{False}) το οποίο μας δίνει την δυνατότητα να αναπαραστήσουμε τις τιμές του $\lambda_+, \lambda_-$ σαν διδιάστατες εικόνες βοηθώντας μας να επιβεβαιώσουμε την αντίστοιχη θεωρία.


Για την συνθήκη (Σ1) χρησιμοποιούμε το μη γραμμικό φίλτρο \foreignlanguage{english}{dilate} και τον ορισμό του ως φίλτρο κινούμενου μεγίστου ώστε να βρούμε τα τοπικά μέγιστα των τιμών του κριτηρίου γωνιότητας που προέκυψαν βάσει της εικόνας που δόθηκε ως όρισμα.

\subsection*{Παρατηρήσεις Αποτελεσμάτων}
\hspace{0.6cm}Στα διαγράμματα αναπαράστασης των ιδιοτιμών του τανυστή (Σχήμα 9) το $\lambda_+$ εμφανίζει μεγάλες τιμές κατά μήκος των ακμών οι οποίες παρουσιάζουν μεγάλη μεταβολή της έντασης σε μία κατεύθυνση που ορίζεται από το $w_+$. Αντίστοιχα το  $\lambda_-$ λαμβάνει μεγάλες τιμές (ταυτόγχρονα με το  $\lambda_+$) στα γωνιακά σημεία δηλαδή εκεί που ακόμα και η ελάχιστη μεταβολή φωτεινότητας (διευθυνσιοδοτείται από το $w_-$) λαμβάνει συγκρητικά μεγάλες τιμές. 

\begin{figure}[H] % Το [h] σημαίνει "here" (εδώ ακριβώς)
    \centering    % Στοίχιση στο κέντρο
    \includegraphics[scale=0.45]{figs_part2/fig8.png}
    %\caption{Παράδειγμα ανιχνευθέντων ακμών για συνεχείς τιμές σ} % Λεζάντα
    \label{fig:results} % Ετικέτα για αναφορά μέσα στο κείμενο (π.χ. Βλέπε Εικόνα \ref{fig:results})
\end{figure}

\begin{figure}[H] % Το [h] σημαίνει "here" (εδώ ακριβώς)
    \centering    % Στοίχιση στο κέντρο
    \includegraphics[scale=0.45]{figs_part2/fig9.png}
    %\caption{Παράδειγμα ανιχνευθέντων ακμών για συνεχείς τιμές σ} % Λεζάντα
    \label{fig:results} % Ετικέτα για αναφορά μέσα στο κείμενο (π.χ. Βλέπε Εικόνα \ref{fig:results})
\end{figure}


\begin{figure}[H] % Το [h] σημαίνει "here" (εδώ ακριβώς)
    \centering    % Στοίχιση στο κέντρο
    \includegraphics[scale=0.45]{figs_part2/fig10.png}
    \caption{Απεικόνιση Ιδιοτιμών} % Λεζάντα
    \label{fig:results} % Ετικέτα για αναφορά μέσα στο κείμενο (π.χ. Βλέπε Εικόνα \ref{fig:results})
\end{figure}


Στην εικόνα \foreignlanguage{english}{solar} μετά την εφαρμογή του αλγορίθμου βλέπουμε ότι έχει εντοπιστεί ένα ικανοποιητικό πλήθος γωνιών στα διαχωριστικά των panel και στις σκιές που αυτά δημιουργουν. Ωστόσο η σημαντική πλειοψηφία των σημείων αυτών βρίκεται σε panel κλίμακας σχετικά μεγάλης (δεδομένου του σ που έχει εφαρμοσθεί) ενώ αγνοούνται οι γωνίες \foreignlanguage{english}{panel} στο βάθος της εικόνας αλλά και οι γωνίες που σχηματίζουν οι εσωτερικές κυψέλες των \foreignlanguage{english}{panels} που αντιστοιχούν σε μικρής κλίμακας αντικείμενα.


Στα \foreignlanguage{english}{blood cells} βλέπουμε τον εντοπισμό γωνιών στα σημεία επαφής των κυττάρων καθώς είναι τα κυριότερα σημεία που μπορούν να προσομοιάσουν γωνίες (τα κύτταρα είναι εν γένει κυκλικά)


Στην τεχνιτή εικόνα \foreignlanguage{english}{edgetest}, χρησιμοποιώντας τις ενδεικτικές παραμέτρους ο αλγόριθος αναγνώρισε όλες τις γωνίες εκτός από τις εξωτρικές γωνίες του γράμματος Τ όπου διαπιστώνεται πολύ μικρότερη διαβάθμιση της φωτεινότητας. Για τον λόγο αυτό επιλέχθηκαν μικρότερες τιμές για τα σ, θ (το οποίο επιβάλλεται λόγω απουσίας θορύβου) ώστε οι γωνίες αυτές να γίνουν ανιχνεύσιμες. 


\begin{figure}[H]
     \centering
     \begin{subfigure}[b]{0.45\textwidth}
         \centering
         \includegraphics[width=\linewidth]{figs_part2/fig11.png}
         \caption{Εικόνα \foreignlanguage{english}{solar.png}}
     \end{subfigure}
     \hfill
     \begin{subfigure}[b]{0.45\textwidth}
         \centering
         \includegraphics[width=\linewidth]{figs_part2/fig12.png}
         \caption{Εικόνα \foreignlanguage{english}{blood\_cells.png}}
     \end{subfigure}
     %\caption{Σύγκριση feature maps.}
\end{figure}

\begin{figure}[H]
     \centering
     \begin{subfigure}[b]{0.45\textwidth}
         \centering
         \includegraphics[width=\linewidth, height=5.5cm]{figs_part2/fig13.png}
         %\caption{Εικόνα \foreignlanguage{english}{solar.png}}
     \end{subfigure}
     \hfill
     \begin{subfigure}[b]{0.45\textwidth}
         \centering
         \includegraphics[width=\linewidth, height=5.5cm]{figs_part2/fig14.png}
         %\caption{Εικόνα \foreignlanguage{english}{blood\_cells.png}}
     \end{subfigure}
     \caption{Ανιχνευθέντες γωνίες στην \foreignlanguage{english}{edgetest.jpg} για διαφορετικές παραμέτρους}
\end{figure}

\subsection*{2.2 Πολυκλιμακωτή Ανίχνευση Γωνιών (\foreignlanguage{english}{Harris-Laplacian})}

\hspace{0.6cm}Η κλασική μέθοδος \foreignlanguage{english}{Harris} είναι αμετάβλητη ως προς την περιστροφή και τη φωτεινότητα, αλλά παρουσιάζει αδυναμία στην αλλαγή κλίμακας. Για την επίλυση αυτού του προβλήματος, υλοποιείται ο αλγόριθμος \foreignlanguage{english}{Harris-Laplacian}, ο οποίος επιτρέπει τον εντοπισμό χαρακτηριστικών σημείων που είναι σταθερά ανεξάρτητα από την κλίμακα εργασίας.

\subsection*{Θεωρητική Επεξήγηση}

\hspace{0.6cm}Στο πρώτο στάδιο εκτελούμε επαναλμβανόμενα των αλγόριθμο εύρεση γωνιών μονής κλίμακας με διαφορετικές τιμές παραμέτρων σ και ρ. Είναι απαραίτητο το σ και το ρ να βρίσκονται σε άμεση συνάρτηση. Η αύξηση του σ απλώνει τα σημεία ενδιαφέροντος, επομένως αν το παράθυρο ολοκλήρωσης κρατηθεί μικρό δεν θα προλάβει να δει όλο το εύρος της μεταβολής. Με αυτόν τον τρόπο, εξασφαλίζουμε ότι αν μια γωνία είναι πολύ μεγάλη για να ανιχνευθεί σε μια λεπτή κλίμακα, θα εντοπιστεί σε μια μεταγενέστερη, πιο εξομαλυμένη κλίμακα όπου η δομή της θα είναι συμβατή με το μέγεθος του παραθύρου ολοκλήρωσης. 

Το δεύτερο στάδιο αφορά την αυτόματη επιλογή της βέλτιστης κλίμακας για κάθε σημείο, χρησιμοποιώντας την κανονικοποιημένη Λαπλασιανή της Γκαουσιανής (\foreignlanguage{english}{LoG}) η οποία είναι μια μετρική αμετάβλητη ως προς την κλίμακα εργασίας. Η χρησιμότητα της εύρεσης τοπικών μεγίστων της \foreignlanguage{english}{LoG} στον χώρο κλίμακας έγγυται στο γεγονός οτι το σημείο ενδιαφέροντος θα προσδιορισθεί στην κλίμακα εκείνη που αντιστοιχεί καλύτερα στις διαστάσεις του στην εικόνα.  


\subsection*{Παρατηρήσεις Αποτελεσμάτων}
\hspace{0.6cm}Για την εικόνα \foreignlanguage{english}{solar} (Σχήμα 12α) παρατηρούμε ότι πλέον κάποιες γωνίες περικυκλώνονται από μεγαλύτερα διακριτικά ανίχνευσης γεγονός που σε συνδιασμό με τις μεγάλες τους διαστάσεις μας οδηγεί στο συμέρασμα ότι ανιχνέυθηκαν σε μεγαλύτεηρ κλίμακα.

Αξιοσημείωτο είναι ακόμα ότι στην εικόνα \foreignlanguage{english}{blood cells} (Σχήμα 12β) οι γωνίες που προέκυπταν από τα σημεία επαφής πλέον δεν ανιχνεύονται (αδύναμο στίγμα στην κλίμακα τους). Αντί για αυτές ως γωνίες πλέον ο αλγόριθμος αναγνωρίζει ολόκληρα τα κύτταρα, τα οποία σε μεγάλες κλίμακες επεξεργασίας μοιάζουν με γωνίες, αφού παρουσιάζουν μεγάλη μεταβολή έντασης ως προς ορθωγώνιες διευθύνσεις (χρωματική αντίθεση \foreignlanguage{english}{cell-background}).

\begin{figure}[H]
     \centering
     \begin{subfigure}[b]{0.45\textwidth}
         \centering
         \includegraphics[width=\linewidth]{figs_part2/fig15.png}
         \caption{Εικόνα \foreignlanguage{english}{solar.png}}
     \end{subfigure}
     \hfill
     \begin{subfigure}[b]{0.45\textwidth}
         \centering
         \includegraphics[width=\linewidth]{figs_part2/fig16.png}
         \caption{Εικόνα \foreignlanguage{english}{blood\_cells.png}}
     \end{subfigure}
     \caption{Πολυκλιμακωτή Ανίχνευση Γωνιών}
\end{figure}

\subsection*{Υλοποίηση σε \foreignlanguage{english}{Python}}
\hspace{0.6cm}Ιδιαίτερη προσοχή δόθηκε στην υλοποίηση του κριτηρίου μεγιστοποίησης της κανονικοποιημένης \foreignlanguage{english}{LoG} και πιο συγκερκιμένα στις ακριανές κλίμακες όπου δεν εντοπίζεται προηγούμενη και επόμενη τιμή. Σε αυτές τις θέσεις έγινε \foreignlanguage{english}{pad} με την τιμή -1 για να μην αλλοιωθεί το αποτέλεσμα.

\subsection*{2.3 Ανίχνευση \foreignlanguage{english}{Blobs}}

\subsection*{Θωρητική Επεξήγηση}
\hspace{0.6cm}Τα ιδιοδιανύσματα του πίνακα \foreignlanguage{english}{Hessian} δίνουν τις κατευθύνσεις της μέγιστης και της ελάχιστης καμπυλότητας και οι ιδιοτιμές την ένταση-μέγεθος της καμπυλότητας. Τα \foreignlanguage{english}{Blobs} είναι περιοχές που παρουσιάζουν έντονες μεταβολές καμπυλότητας και προς τις δύο κατευθύνσεις. Βάση της ιδιότητας αυτής των \foreignlanguage{english}{Blobs}, το κριτήριο ανίχνευσής τους ορίζεται ως η ορίζουσα του πίνακα \foreignlanguage{english}{Hessian}, δηλαδή το γινόμενο των ιδιοτιμών του. Όταν και οι δύο ιδιοτιμές είναι μεγάλες τότε το κριτήριο λαμβάνει μεγάλη τιμή υπερβαίνοντας το προκαθορισμένο κατώφλι ανίχνευσης.

\subsection*{Παρατηρήσεις Αποτελεσμάτων}
\hspace{0.6cm}Γνωρίζουμε ότι τα \foreignlanguage{english}{Blobs} είναι περιοχές των εικόνων που είναι σημαντικά πιο φωτεινές ή πιο σκοτεινές από το περιβάλλον γύρω τους, δηλαδή προσομοιάζουν κηλίδες.

Στην εικόνα \foreignlanguage{english}{Blood Cells} (Σχήμα 14), όπου η δομή των κυττάρων συγκλίνει στον ορισμό των \foreignlanguage{english}{Blobs}, το αποτέλεσμα της ανίχνευσης είναι το διαισθητικά αναμενόμενο. Πιο συγκεκριμένα βλέπουμε ολόκληρα κύτταρα να έχουν περικυκλωθεί απο ανιχνευτικούς κύκλους, ωστόσο για να συμβεί αυτό πρέπει να αυξήσουμε το σ, ώστε το μέγεθος του παραθύρου αναζήτησης να αντιστοιχεί στην ακτίνα των κυττάρων. Αντίθετα για μικρό σ ο αλγόριθμος δεν λαμβάνει υπόψη ολόκληρο το κύτταρο αλλά εντοπίζει μικρές ανωμαλίες στην επιφάνειά του. Δεν πρέπει να αποτελεί έκλπηξη η ανίχνευση λευκών περιοχών του \foreignlanguage{english}{background} που σχηματίζονται σαν ````λίμνες'''' σε σημεία σύγκλισης πολλαπλών κυττάρων.

Στην εικόνα \foreignlanguage{english}{Solar} (Σχήμα 13), όπου δεν εμφανίζονται δομές με την ````παραδοσιακή'''' έννοια των \foreignlanguage{english}{Blobs}, έχουμε ανίχνευση κάποιων συμπαγών σχηματισμών στο \foreignlanguage{english}{background} και επιπλέον, ιδιαίτερα στην μικρή κλίμακα, ανίχνευση των σημείων συνάντησης των γωνιών των \foreignlanguage{english}{Panel}. Η ανίχνευση των τελευταίων αυτών σημείων μπορεί παραπλανητικά να θεωρηθεί ανίχνευση γωνιών καθώς σε πολύ μικρή κλίμακα μπορεί να ειδωθεί ως ένα \foreignlanguage{english}{Blob}.

\begin{figure}[H]
     \centering
     \begin{subfigure}[b]{0.45\textwidth}
         \centering
         \includegraphics[width=\linewidth]{figs_part2/fig17.png}
         %\caption{Εικόνα \foreignlanguage{english}{solar.png} (σ = )}
     \end{subfigure}
     \hfill
     \begin{subfigure}[b]{0.45\textwidth}
         \centering
         \includegraphics[width=\linewidth]{figs_part2/fig18.png}
         %\caption{Εικόνα \foreignlanguage{english}{blood\_cells.png}}
     \end{subfigure}
     \caption{Σύγκριση Ανίχνευσης \foreignlanguage{english}{Blobs} στην εικόνα \foreignlanguage{english}{solar.png} για σ = 2 και σ = 6}
\end{figure}


\begin{figure}[H]
     \centering
     \begin{subfigure}[b]{0.45\textwidth}
         \centering
         \includegraphics[width=\linewidth]{figs_part2/fig19.png}
         %\caption{Εικόνα \foreignlanguage{english}{solar.png} (σ = )}
     \end{subfigure}
     \hfill
     \begin{subfigure}[b]{0.45\textwidth}
         \centering
         \includegraphics[width=\linewidth]{figs_part2/fig20.png}
         %\caption{Εικόνα \foreignlanguage{english}{blood\_cells.png}}
     \end{subfigure}
     \caption{Σύγκριση Ανίχνευσης \foreignlanguage{english}{Blobs} στην εικόνα \foreignlanguage{english}{blood\_cells.png} για σ = 2 και σ = 6}
\end{figure}


\subsection*{2.4. Πολυκλιμακωτή Ανίχνευση \foreignlanguage{english}{Blobs}}
\hspace{0.6cm}Ο κλασσικός ανιχνευτής Hessian επεκτείνεται εδώ υιοθετώντας τις μεθόδους τις πολυκλιμακωτής ανάλυσης.

\subsection*{Παρατηρήσεις Αποτελεσμάτων}
\hspace{0.6cm}Στην εικόνα \foreignlanguage{english}{Blood Cells} (Σχήμα 15β) τόσο μικρά όσο και μεγάλα κύτταρα περιβάλλονται από μικρούς και μεγάλους κύκλους αντίστοιχα, γεγονός το οποίο υποδηλώνει την διαφορετική κλίμακα ανίχνευσής τους. Οι εμφολευμνένοι κύκλοι που παρατηρούμε οφείλονται στις ανομοιομορφίες στην επιφάνεια μεγαλύτερων κυττάρων. 

Παρόμοια και στην εικόνα \foreignlanguage{english}{Solar} (Σχήμα 15α) λεπτά χαρακτηριστικά επισημαίνονται με μικρούς κύκλους ενώ πιο ευμεγέθη με κύκλους μεγαλύτερης ακτίνας. Σε αυτήν την εικόνα η πιο χρήσιμη πληροφορία φαίνεται να εντοπίζεται στους μεγαλύτερους κύκλους. 

Με την πολυκλιμακωτή ανάλυση αυτή τα αποτελέσματα του κλασσικού αλγόριθμου \foreignlanguage{english}{Hessian} σε πολλές κλίμακες υπερτίθενται σε μία εικόνα παρέχοντας μια συμπυκνωμένη πληροφορία για τα χαρακτηριστικά πολλών διαφορετικών κλιμάκων.

\begin{figure}[H]
     \centering
     \begin{subfigure}[b]{0.45\textwidth}
         \centering
         \includegraphics[width=\linewidth]{figs_part2/fig21.png}
         \caption{Εικόνα \foreignlanguage{english}{solar.png}}
     \end{subfigure}
     \hfill
     \begin{subfigure}[b]{0.45\textwidth}
         \centering
         \includegraphics[width=\linewidth]{figs_part2/fig22.png}
         \caption{Εικόνα \foreignlanguage{english}{blood\_cells.png}}
     \end{subfigure}
     \caption{Πολυκλιμακωτή Ανίχνευση  \foreignlanguage{english}{Blobs}}
\end{figure}

\subsection*{Χρήσιμα Συμπεράσματα Μέρους 2}
\hspace{0.6cm}Από την ενασχόληση μας τόσο με τον ανιχνευτή γωνιών \foreignlanguage{english}{Harris} όσο και με τον ανιχνευτηή \foreignlanguage{english}{Blobs Hessian}, αλλά και με τις πολυκλιμακωτές τους εκδοχές, αντιληφθήκαμε ότι η εικόνα \foreignlanguage{english}{Solar} προσφέρεται για εντοπισμό γωνιών ενώ η εικόνα \foreignlanguage{english}{Blood Cells} προορίζεται για την εύρεση \foreignlanguage{english}{Blobs}. Το γεγονός ότι λάβαμε επισημειώσεις και για τις δύο εικόνες στους δύο αλγορίθμους αποδεικνύει την ευκολία με την οποία οι αλγόριθμοι μπορεί να παρανοήσουν, ανάλογα με την κλίμακα επεξεργασίας, την πρωταρχική έννοια των χαρακτηριστικών (σύγχυση γωνίας-\foreignlanguage{english}{blob}). Για παράδειγμα ένα στρογγυλό αντικείμενο μπορεί αν ανιχνευθεί ως γωνία σε μεγάλες κλίμακες και αντίστοιχα μια γωνία ως \foreignlanguage{english}{blob} σε μικρές κλίμακες, πράγμα που αποκαλύπτει ένα ευάλωτο σημείο στην χρήση μαθηματικών εργαλείων από τους αλγορίθμους. 

Το πρόβλημα αυτό αναβαθμίζει την δυσκολία εύρεσης των κατάλληλων παραμέτρων των παραπάνω μεθόδων που θα καταστήσουν τα αποτελέσματά τους όσο το δυνατόν εγκυρότερα και πιο αξιόπιστα. Οι ιδανικές παράμετροι, μάλιστα, μεταβάλλονται δραστικά από φωτογραφία σε φωτογραφία ανάλογα με τις προθέσεις μας και τα εγγενή χαρακτηριστικά της.

\section*{Μέρος 3: Κατηγοριοποίηση Εικόνων με Χρήση Τοπικών Περιγραφητών
και Χαρακτηριστικών από Βαθύ Νευρωνικό Δίκτυο}

\section*{Βήμα 3.1}

\hspace{0.6cm}Η εξαγωγή περιγραφητών συμπυκνώνει την οπτική πληροφορία σε διανύσματα για την αναγνώριση προτύπων. Τόσο ο τοπικός \foreignlanguage{english}{HOG (Local HOG)}, μια παραλλαγή του κλασσικού \foreignlanguage{english}{HOG} όσο και ο \foreignlanguage{english}{SURF} εστιάζουν αποκλειστικά γύρω από προ-ανιχνευμένα σημεία ενδιαφέροντος, προσαρμόζοντας το μέγεθος του τοπικού παραθύρου ανάλυσης στην κλίμακα $\sigma$ στην οποία ανιχνεύθηκε το εκάστοτε σημείο.

\subsection*{Τοπικός \foreignlanguage{english}{HOG}}
\hspace{0.6cm}Σε αυτή τη μορφή, ο \foreignlanguage{english}{HOG} δεν εξετάζει το σύνολο της εικόνας, αντίθετα, εξάγει τον εκάστοτε περιγραφητή από ένα τοπικό παράθυρο (\foreignlanguage{english}{patch}) με κέντρο το σημείο ενδιαφέροντος, εφαρμόζοντας την χωρική διαμέριση σε (επικαλυπτόμενα) κελιά αποκλειστικά στο παράθυρο αυτό. 

\begin{itemize}
    \item \textbf{Εντοπισμός Άκαμπτων Αντικειμένων:} Ο τοπικός \foreignlanguage{english}{HOG} δεν περιστρέφει το παράθυρο βάσει κάποιας τοπικής κύριας κατεύθυνσης. Επομένως, διατηρεί την απόλυτη χωρική κατανομή των τοπικών κλίσεων. Επισπρόσθετα η πυκνή διαμέριση σε επικαλυπτόμενα κελιά εξασφαλίζει την συνέχεια των περιγραμμάτων των αντικειμένων διατηρώντας την συσχετιση των γειτονικών κελιών. Πρακτικά, αυτός ο συνδυασμός τον καθιστά εξαιρετικό στην ταύτιση άκαμπτων γεωμετρικών δομών (π.χ. τεχνητά αντικείμενα, κτίρια, πάνελ), όπου η απόλυτη γεωμετρία και η γωνία των ακμών είναι σταθερές και δομικά χαρακτηριστικές, με την προϋπόθεση ότι δεν υφίσταται έντονη περιστροφή της κάμερας.
    \item \textbf{Κανονικοποίηση:} Η χρήση της $\ell_1$ ή $\ell_2$ νόρμας σε επίπεδο τοπικού μπλοκ καθιστά τον περιγραφητή εξαιρετικά εύρωστο σε τοπικές μεταβολές του φωτισμού, καθώς οι κλίσεις αξιολογούνται σχετικά με τη γειτονιά τους και όχι με βάση την απόλυτη φωτεινότητα.
\end{itemize}

\subsection*{\foreignlanguage{english}{SURF}}
\hspace{0.6cm}Ο \foreignlanguage{english}{SURF} κατασκευάζει την περιγραφή αξιοποιώντας αποκρίσεις φίλτρων διακριτής προσέγγισης της παραγώγου (\foreignlanguage{english}{Haar Wavelets}) στην τοπική γειτονιά του σημείου, με έμφαση στην υπολογιστική ταχύτητα και τη μη μεταβολή της γεωμετρίας.

\begin{itemize}
    \item \textbf{Ανάγκη για Διακριτή Τοπική Υφή:} Ο υπολογισμός της κύριας κατεύθυνσης, καθώς και των αθροισμάτων του ίδιου του περιγραφητή, βασίζεται αποκλειστικά στις αποκρίσεις των τοπικών ακμών (διαφορές φωτεινότητας). Συνεπώς, η μέθοδος απαιτεί έντονη τοπική υφή ή υψηλή αντίθεση. Σε ομοιογενείς περιοχές, οι κλίσεις είναι πρακτικά μηδενικές, καθιστώντας τον προσανατολισμό ασταθή και τον περιγραφητή μη-πληροφοριακό, γεγονός που οδηγεί σε εσφαλμένες ταυτίσεις.
    \item \textbf{Κύρια Κατεύθυνση και Ανεξαρτησία Περιστροφής:} Εφόσον υφίσταται επαρκής υφή, ο \foreignlanguage{english}{SURF} εστιάζει σε μια γειτονιά του σημείου ενδιαφέροντος ανάλογη της κλίμακας ανίχνυεσης και την περιστρέφει βάσει της επικρατούσας κατεύθυνσης. Αυτό εξασφαλίζει πλήρη ανεξαρτησία περιστροφής, επιτρέποντας την ταύτιση μοτίβων υπό οποιαδήποτε γωνία λήψης.
    \item \textbf{Γκαουσιανά Βάρη:} Η απόδοση μεγαλύτερου βάρους στις κλίσεις κοντά στο κέντρο του παραθύρου ελαχιστοποιεί την επίδραση της περιφέρειας.
\end{itemize}

\subsection*{Συγκριτική Σύνοψη}
\hspace{0.6cm}Ο \textbf{Τοπικός \foreignlanguage{english}{HOG}} πλεονεκτεί στην ταύτιση άκαμπτων δομών με σταθερό προσανατολισμό, κωδικοποιώντας την απόλυτη χωρική διάταξη των ακμών με μεγάλη ακρίβεια, θυσιάζοντας όμως την ανεξαρτησία περιστροφής. Αντίθετα, ο \textbf{\foreignlanguage{english}{SURF}} είναι η βέλτιστη επιλογή για αντιστοίχιση εικόνων υπό ελεύθερες συνθήκες λήψης (\foreignlanguage{english}{scale/rotation invariance}), με την αυστηρή προϋπόθεση ύπαρξης πλούσιας και διακριτής τοπικής υφής.

\subsection*{Παρατηρήσεις Αποτελεσμάτων}
\hspace{0.6cm}Με βάση την μετρική \foreignlanguage{english}{Accuracy} (Σχήμα 16) είναι φανερό ότι ο εξαγωγέας περιγραφητών \foreignlanguage{english}{SURF} υπερέχει του \foreignlanguage{english}{HOG} σε όλους τους συνδυασμούς. Αυτό μοιάζει να οφείλεται στην ανοχή που επιδεικνύει σε τυχόν περιστροφές της εικόνας, πράγμα που συμβαίνει σε ορισμένες εικόνες του δοθέντος \foreignlanguage{english}{dataset}. 

Με βάση την ίδια μετρική οι ανιχνευτές \foreignlanguage{english}{Hessian-Laplacian} (ανιχνευτής \foreignlanguage{english}{Blobs}) και \foreignlanguage{english}{Harris-Laplace} (ανιχνευτής γωνιών) φαίνεται να δίνουν παρόμοια αποτελέσματα για κοινό περιγραφητή. Το ποιός από τους δύο υπερέχει προκύπτει από το είδος των εικόνων στο οποίο εφαρμόζονται (μετά από τυχαία διάσπαση σε \foreignlanguage{english}{train} και \foreignlanguage{english}{test}), δηλαδή από το αν τα ουσιαστικά χαρακτηριστικά των εικόνων είναι γωνίες ή \foreignlanguage{english}{blobs}.

\begin{figure}[H] % Το [h] σημαίνει "here" (εδώ ακριβώς)
    \centering    % Στοίχιση στο κέντρο
    \includegraphics[scale=0.45]{figs_part3/fig23.png}
    \caption{\foreignlanguage{english}{Accuracy} και \foreignlanguage{english}{Standard Deviation} για διαφορετικούς συνδυασμούς Ανιχνευτή-Περιγραφητή} % Λεζάντα
    \label{fig:results} % Ετικέτα για αναφορά μέσα στο κείμενο (π.χ. Βλέπε Εικόνα \ref{fig:results})
\end{figure}

Η ανάλυση των \foreignlanguage{english}{confusion matrices} (Σχήμα 17) αναδεικνύει χρήσιμα συμπεράσματα. Ο \foreignlanguage{english}{HOG} αποτυγχάνει παταγωδώς στην αναγνώριση ανθρώπων πράγμα που αντιβαίνει τα προσδοκόμενα αποτελέσματα βάσει της θεωρίας του κλασσικού \foreignlanguage{english}{HOG}. Αυτό γιατί, η κλασσική εκδοχή του σαρώνει όλη την εικόνα και βλέπει το κεφάλι, τους ώμους και τα πόδια ως μια ενιαία σιλουέτα. Αυτή η σιλουέτα είναι το ισχυρότερο χαρακτηριστικό ενός ανθρώπου. Αντίθετα, ο \foreignlanguage{english}{Local HOG} βασίζεται σε συγκεκριμένα σημεία ενδιφέροντος μικρότερης κλίμακας που του παρέχουν οι ανιχνευτές π.χ. έναν αγκώνα ή μια πτυχή του ρούχου. Ο \foreignlanguage{english}{HOG} μετά περιγράφει μόνο αυτό το μικρό τετραγωνάκι (\foreignlanguage{english}{patch}). Σε αυτή την κλίμακα, ο \foreignlanguage{english}{HOG} δεν βλέπει ````άνθρωπο'''', βλέπει απλώς τυχαίες γραμμές και κλίσεις. Ωστόσο ο \foreignlanguage{english}{HOG} προσεγγίζει εξαιρετικα την επίδοση του \foreignlanguage{english}{SURF} στην αναγώριση ποδηλάτων. Αυτό γιατί στα ποδήλατα ακόμα και ένα μικρό \foreignlanguage{english}{patch} είναι χαρακτηριστικό. Αν το \foreignlanguage{english}{patch} πέσει πάνω στις ακτίνες ή στο τιμόνι, οι κλίσεις είναι πολύ συγκεκριμένες (γεωμετρικές). Ο συνδυασμός πολλών τέτοιων ````γεωμετρικών'''' \foreignlanguage{english}{patches} στο \foreignlanguage{english}{Bag of Words} αρκεί για να αναγνωριστεί το ποδήλατο.

Προτείνουμε για την κατηγοριοποίηση των εικόνων τον συνδυασμό \foreignlanguage{english}{Hessian-Laplace + SURF}. Βάσει των αποτελεσμάτων ο συνδυασμός αυτός έχει τόσο το μεγαλύτερο \foreignlanguage{english}{Accuracy} όσο και το μικρότερο \foreignlanguage{english}{Standard Deviation} (\foreignlanguage{english}{Accuracy} = 72.828\%, \foreignlanguage{english}{Standard Deviation} = 2.602\%).

\begin{figure}[H] % Το [h] σημαίνει "here" (εδώ ακριβώς)
    \centering    % Στοίχιση στο κέντρο
    \includegraphics[scale=0.4]{figs_part3/fig24.png}
    \caption{\foreignlanguage{english}{Confusion Matrices} για διαφορετικούς συνδυασμούς Ανιχνευτή-Περιγραφητή} % Λεζάντα
    \label{fig:results} % Ετικέτα για αναφορά μέσα στο κείμενο (π.χ. Βλέπε Εικόνα \ref{fig:results})
\end{figure}

\subsection*{Πειραματισμός-Οπτικοποίηση Σημείων Ενδιαφέροντος εικόνων του \foreignlanguage{english}{Dataset}}
\hspace{0.6cm}Στην προσπάθεια να κατανοήσουμε καλύτερα και να συγκεκριμενοποιήσουμε τα σημεία ενδιαφέροντος για τις συγκεκριμένες κατηγορίες εικόνων (άνρθωπος, αυτοκίνητο, ποδήλατο) οπτικοποιούμε τα αποτελέσματα του Πολυκλιμακωτού Ανιχνευτή  \foreignlanguage{english}{Blobs} σε 3 τυχαίες φωτογραφίες του συνόλου δεδομένων. 



\begin{figure}[H]
     \centering
     \begin{subfigure}[b]{0.45\textwidth}
         \centering
         \includegraphics[width=\linewidth]{figs_part3/fig27.png}
         \caption{Εικόνα \foreignlanguage{english}{bike\_100.png}}
     \end{subfigure}
     \hfill
     \begin{subfigure}[b]{0.45\textwidth}
         \centering
         \includegraphics[width=\linewidth]{figs_part3/fig26.png}
         \caption{Εικόνα \foreignlanguage{english}{carsgraz\_170.png}}
     \end{subfigure}
     \hfill
     \centering
     \begin{subfigure}[b]{0.4\textwidth}
         \centering
         \includegraphics[width=\linewidth]{figs_part3/fig25.png}
         \caption{Εικόνα \foreignlanguage{english}{person\_045.png}}
     \end{subfigure}
     \caption{Ανίχνευση \foreignlanguage{english}{Blobs} σε εικόνες του \foreignlanguage{english}{Dataset}}
\end{figure}


\section*{Βήμα 3.2}

\hspace{0.6cm}Η χρήση Συνελικτικών Νευρωνικών Δικτύων παρακάμπτει την δημιουργία διανυσμάτων χαρακτηριστικών από ````χειροποίητες'''' (αποστασιοποιημένες από το \foreignlanguage{english}{dataset}) μεθόδους και εκμεταλλεύεται την προ-εκπαίδευση σε ένα μεγάλο σύνολο εικόνων ώστε να εξάγει τα επιθυμητά διανύσματα (αναπαραστάσεις που προκύπτουν άμεσα από τα δεδομένα).

\subsection*{Ιεραρχική Εξαγωγή Χαρακτηριστικών στα \foreignlanguage{english}{CNNs}}
\hspace{0.6cm}Ο μετασχηματισμός της εισόδου $X$ σε ένα πολυδιάστατο διάνυσμα χαρακτηριστικών συντελείται μέσω αλλεπάλληλων μη-γραμμικών απεικονίσεων $x^{(l)} = f(W^{(l)} * x^{(l-1)} + b^{(l)})$, όπου $*$ δηλώνει τη συνέλιξη και $f(\cdot)$ μια μη-γραμμική συνάρτηση ενεργοποίησης. 

\begin{itemize}
    \item \textbf{Αρχικά Επίπεδα (\foreignlanguage{english}{Low-level Features)}:} Οι πυρήνες των πρώτων επιπέδων λειτουργούν ως τοπικοί ανιχνευτές συχνοτήτων και κατευθύνσεων. Λόγω του μικρού τοπικού πεδίου λήψης, κωδικοποιούν πρωταρχικές δομές, όπως ακμές και γωνίες.
    \item \textbf{Ενδιάμεσα Επίπεδα (\foreignlanguage{english}{Mid-level Features)}:} Υπολογίζοντας μη-γραμμικούς συνδυασμούς των αρχικών αποκρίσεων, τα επίπεδα αυτά συγκροτούν σύνθετα μοτίβα και τοπικές υφές. Κάθε πυρήνας συνέλιξης στα επίπεδα αυτά 'κοιτάει' συνδυασμένη πληροφορία από ολοένα και μεγαλύτερο εύρος της αρχικής εικόνας.
    \item \textbf{Τελικά Επίπεδα (\foreignlanguage{english}{High-level Semantics}):} Το πεδίο λήψης καλύπτει πλέον το μεγαλύτερο μέρος της εικόνας. Οι ενεργοποιήσεις κωδικοποιούν αφηρημένες έννοιες (π.χ. `````τροχός'''', ````πρόσωπο'''') επιδεικνύοντας αμεταβλητότητα σε μικρές χωρικές μεταβολές.
\end{itemize}

Τα αποτελέσματα της οπτικοποίησης των \foreignlanguage{english}{layers} 1, 3, 6 και 12 μιας τυχαίας εικόνας ανθρώπου παρατίθεται στο Σχήμα 19. Τα πρώτα επίπεδα (\foreignlanguage{english}{layer} 1 και  \foreignlanguage{english}{layer2}), εμφανίζουν αρκετά υψηλή ανάλυση και αποτύπωση της πραγματικότητας και εντοπίζουν βασικά σημεία ενδιαφέροντος, όπως ακμές, γωνίες και αλλαγές φωτισμού (περίγραμμα ανθρώπινης σιλουέτας και βασικών χαρακτηριστικών προσώπου και ρούχων, σκελετός ποδηλάτου και σχήμα σέλας). Οι συνελίξεις, οι μη γραμμικές ενεργοποιήσεις και τα \foreignlanguage{english}{pooling} επίπεδα που μεσολαβούν αλλοιώνουν την ανάλυση της εικόνας. Τελικά προκύπτουν πολύ αφηρημένες εικόνες μικρής ανάλυσης που δεν θυμίζουν επουδενί τις αρχικές αλλά αποτελούν μια γενικευμένη αποτύπωση με υψηλές τιμές φωτεινότητας στα σημεία που εντοπίζεται το αντικείμενο ενδιαφέροντος (άνθρωπος, ποδήλατο και αμάξι).

\begin{figure}[H] % Το [h] σημαίνει "here" (εδώ ακριβώς)
    \centering    % Στοίχιση στο κέντρο
    \includegraphics[scale=0.4]{figs_part3/fig28.png}
    \caption{Οπτικοποίση ενδιάμεσων \foreignlanguage{english}{layers} του \foreignlanguage{english}{CNN}} % Λεζάντα
    \label{fig:results} % Ετικέτα για αναφορά μέσα στο κείμενο (π.χ. Βλέπε Εικόνα \ref{fig:results})
\end{figure}

\subsection*{Η Σημασία της Κανονικοποίησης}
\hspace{0.6cm}Η προεπεξεργασία της εισόδου μέσω του μετασχηματισμού $X_{norm} = \frac{X - \mu}{\sigma}$ (όπου $\mu, \sigma$ οι παράμετροι του συνόλου εκπαίδευσης) είναι επιβεβλημένη προϋπόθεση. Κατά την εκπαίδευση του νευρωνικού δικτύου τα δεδομένα έχουν κανονικοποιηθεί σε επίπεδο \foreignlanguage{english}{batch} με σκοπό την ταχύτερη σύγκλιση του αλγορίθμου εκμάθησης. Συνεπώς κατά την πρόβλεψη τα νέα δείγματα οφείλουν να υπόκεινται ακρβώς στον ίδιο μετασχηματισμό με αυτόν που γνωρίζουν τα βελτιστοποιημένα βάρη του δικτύου. Απόκλιση από αυτές τις τιμές προκαλεί αυθαίρετες και εσφαλμένες ενεργοποιήσεις των μη-γραμμικών συναρτήσεων.


\subsection*{Υπεροχή έναντι Παραδοσιακών Μεθόδων}
\hspace{0.6cm}Τα \foreignlanguage{english}{CNNs} υπερτερούν σημαντικά έναντι των μεθόδων \foreignlanguage{english}{Bag-of-Visual-Words}. Αυτό καθώς αντί για χειροκίνητα σχεδιασμένους περιγραφητές, το δίκτυο, μέσω της εκπαίδευσής του, μαθαίνει τους βέλτιστους πυρήνες συνέλιξης για το συγκεκριμένο πρόβλημα. Επιπρόσθετα, το \foreignlanguage{english}{BoVW} καταστρέφει τη χωρική πληροφορία μέσω της κβαντοποίησης (ιστόγραμμα), ενώ το \foreignlanguage{english}{CNN} διατηρεί τη χωρική συσχέτιση των χαρακτηριστικών μέσω των \foreignlanguage{english}{feature maps}. Τέλος το \foreignlanguage{english}{CNN} μετασχήματίζει τον αρχικό χώρο περιγραφής των εικόνων σε έναν άλλο μεγαλύτερης διάστασης που ενδέχεται να καθιστά το πρόβλημα γραμμικά διαχωρίσιμο και κατ' επέκταση ευεπίλυτο από έναν απλό \foreignlanguage{english}{SVM} ταξινομητή.


δίνω accuracy και std των accuracy δινω conf matrices

Αυτή η υπεροχή επιβεβαιώνεται εμφατικά από τα πειράματα. Βλέπουμε άυξηση στην μετρική \foreignlanguage{english}{accuracy} 20 ποσοστιαίων μονάδων (\foreignlanguage{english}{Accuracy} = 95.31\%) (μειωση std ενδεχομένως). Το αδύναμο σημείο της αναγνώρισης του ανθρώπου στους τοπικούς περιγραφητές διορθώνεται ικανοποιητικά (από σχεδόν τυχαία ταξινόμηση πάμε σε ποσοστά επιτυχίας περίπου 90\%) αλλά δεν παύει να παραμένει το χαμηλότερο ποσοστό ανάμεσα στις 3 κατηγορίες όπως φαίνεται από το Confusion Matrix του Σχήματος

\begin{figure}[H] % Το [h] σημαίνει "here" (εδώ ακριβώς)
    \centering    % Στοίχιση στο κέντρο
    \includegraphics[scale=0.7]{figs_part3/fig29.png}
    \caption{\foreignlanguage{english}{Confusion Matrix Clean Dataset - CNN}} % Λεζάντα
    \label{fig:results} % Ετικέτα για αναφορά μέσα στο κείμενο (π.χ. Βλέπε Εικόνα \ref{fig:results})
\end{figure}

\subsection*{Θόρυβος και Περιστροφή - \foreignlanguage{english}{CNN}}

Η εισαγωγή παραμορφώσεων δοκιμάζει την ανθεκτικότητα των εξαγόμενων χαρακτηριστικών: 

\begin{itemize}
    \item \textbf{Γκαουσιανός Θόρυβος}: Τα \foreignlanguage{english}{CNNs} παρουσιάζουν εγγενή ανθεκτικότητα στον λευκό θόρυβο. Οι πράξεις υποδειγματοληψίας (\foreignlanguage{english}{Max/Average Pooling}) λειτουργούν πρακτικά ως βαθυπερατά φίλτρα, καταστέλλοντας τις υψίσυχνες χωρικές διακυμάνσεις του θορύβου. Τα τελευταία επιπίπεδα άλλωστε εξάγουν πληροφορία εξαιρετικά γενικευμνένη που η επιρροή των τοπικά εντοπισμένων μεταβολων έχει ελαχιστοποιηθεί. Ωστόσο, ακραίος θόρυβος αποδομεί τις ακμές στα αρχικά επίπεδα, προκαλώντας διάδοση σφάλματος στον τελικό περιγραφητή.
    \item \textbf{Περιστροφή:} Οι κλασικές πράξεις συνέλιξης είναι ισοδύναμες ως προς τη μετατόπιση, αλλά όχι αμετάβλητες ως προς την περιστροφή. Αναμένεται μερική διαφοροποίηση των εξαγώμενων διανυσμάτων χαρακτηριστικών συγκριτικά με αυτά των μη περιστραμμένων εικόνων. Η διαφοροποίηση αυτή όμως είναι περιορισμένη καθώς η μορφή που παίρνει το τελικό διάνυσμα στα ανώτερα επίπεδα διαμορφώνεται από την ύπαρξη μιας αφηρημένης δομής (π.χ σέλα), που το δίκτυο θεωρεί σημαντική, και όχι από τον προσανατολισμό της. Άρα τα τελικά διανύσματα που προκύπτουν από την ίδια κλάση κατλαμβάνουν συμπαγή περιοχή του χώρου και παραμένουν εύκολα διαχωρίσιμα από τον τελικό ταξινομητή. Αναμένεται σημαντική πτώση απόδοσης, καθώς ένα περιστραμμένο αντικείμενο διεγείρει εντελώς διαφορετικούς πυρήνες στα αρχικά επίπεδα. Εδώ, παραδοσιακές μέθοδοι όπως ο \foreignlanguage{english}{SURF} (λόγω της περιστροφικής ευθυγράμμισης μέσω κυρίαρχης κατεύθυνσης) διαθέτουν θεωρητικό πλεονέκτημα αμεταβλητότητας, σε αντίθεση με ένα τυπικό \foreignlanguage{english}{CNN} που δεν έχει εκπαιδευτεί με εκτεταμένη επαύξηση δεδομένων σε περιστροφές.
\end{itemize}


Πράγματι τα πειραματικά δεδομένα (\foreignlanguage{english}{accuracy, confusion matrices}) των Σχημάτων 21 και 22 οπτικοποιούν την ανοχή του μοντέλου στον θόρυβο και στην περιστροφή. Παρόλα αυτά η παρουσία των παραμορφώσεων αυτών δεν μπορεί να αφήσει το αποτέλεσμα ανεπηρέαστο παρατηρώντας έστω και μικρή αλλά υπαρκτή πτώση στις μετρικές. Πιο συγκεκριμένα παρουσιάζεται πτώση τάξης του 2\% στο \foreignlanguage{english}{mean accuracy} ενώ χαρακτηριστική είναι η πτώση του \foreignlanguage{english}{recall} 4\% στην ταξινόμηση των ανθρώπων όταν επενεργεί συνδυασμός παραμόρφωσης περιστροφής και θορύβου. Σε μια απόπειρα δικαιολόγησης της τελευταίας παρατήρησης, διαπιστώνουμε ότι ο άνθρωπος δεν διαθέται πολλά συμμετρικά χαρακτηριστικά (ευάλωτος σε περιστροφές) ενώ η μορφή του καταλαμβάνει μια μεγάλη επιφάνεια της αρχικής εικόνας (ευάλωτος στην συσσώρευση θορύβου). Αντίθετα στοιχεία του αυτοκινήτου όπως οι ρόδες και τα παράθυρα του αποδίδουν μια συμμετρική δομή ενώ το ποδήλατο είναι εμφανώς εντοπισμένο στον χώρο (λεπτός σκελετός) με πρόσθετα στοιχεία συμμετρίας, τους τροχούς του.

\begin{figure}[H] % Το [h] σημαίνει "here" (εδώ ακριβώς)
    \centering    % Στοίχιση στο κέντρο
    \includegraphics[scale=0.5]{figs_part3/fig30.png}
    \caption{\foreignlanguage{english}{Accuracies Distorted Datasets - CNN}} % Λεζάντα
    \label{fig:results} % Ετικέτα για αναφορά μέσα στο κείμενο (π.χ. Βλέπε Εικόνα \ref{fig:results})
\end{figure}

\begin{figure}[H] % Το [h] σημαίνει "here" (εδώ ακριβώς)
    %\centering    % Στοίχιση στο κέντρο
    \includegraphics[scale=0.25]{figs_part3/fig31.png}
    \caption{\foreignlanguage{english}{Confusion Matrices Distorted Datasets - CNN}} % Λεζάντα
    \label{fig:results} % Ετικέτα για αναφορά μέσα στο κείμενο (π.χ. Βλέπε Εικόνα \ref{fig:results})
\end{figure}



\subsection*{Θόρυβος και Περιστροφή - Παραδοσιακοί Μέθοδοι}
\hspace{0.6cm}Οι παραδοσιακές μέθοδοι εξαγωγής χαρακτηριστικών εμφανίζουν ανεπαίσθητη διαφορά της απόδοσης με την προσθήκη θορύβου στο σύνολο δεδομένων. Η διαφορά αυτή πιθανώς να ήταν πιο διευρυμένη με προσθήκη θορύβου μεγαλύτερης τυπικής απόκλισης. Ακόμα και στις μικρές τιμές θορύβου, ωστόσο, ο συνδυασμός \foreignlanguage{english}{Hessian-Laplace + HOG} παρουσίασε πτώση 6\% στο \foreignlanguage{english}{accuracy}. Η μικρή αλλά υπαρκτή διαφοροποίηση πιθανώς μπορεί αν διακιολογηθεί από το γεγονός ότι τόσο οι ανιχνευτές σημείων ενδιαφέροντος όσο και οι περιγραφητές βασίζονται στον υπολογισμό παραγώγων, διαδικασία η οποία επηρεάζεται έντονα από την παρουσία θορύβου, παρά τις όποιες προσπάθειες ομαλοποίησης με γκαουσιανά φίλτρα. Ως προς την ανίχνευση ο θόρυβος δημιουργεί ````ψευδή'''' σημεία ενδιαφέροντος ή/και αποκρύπτει τα πραγματικά σημεία ενδιαφέροντος.

Με την περιστροφή του συνόλου δεδομένων παρατηρείται η θεμελιώδης διαφορά στη σχεδίαση των δύο περιγραφητών. Ο \foreignlanguage{english}{SURF} έχει σχεδιαστεί για να είναι αναλλοίωτος στην περιστροφή, συνεπώς το \foreignlanguage{english}{accuracy} ναι μεν εμφανίζει πτώση (της τάξης του 5\%) αλλά πιο συγκρατημένη απο αυτήν του \foreignlanguage{english}{HOG} (της τάξης του 9\%) όπως φαίνεται στο Σχήμα 23. Η υλοποίηση του \foreignlanguage{english}{HOG} δεν περιλαμβάνει στάδιο κανονικοποίησης προσανατολισμού. Αν η εικόνα περιστραφεί κατά μια γωνία θ, οι κλίσεις των παραγώγων θα αλλάξουν κατά την ίδια γωνία, αλλάζοντας τελείως τα \foreignlanguage{english}{bins} του ιστογράμματος. Το \foreignlanguage{english}{Bag of Words} θα αποτύχει να αντιστοιχίσει το ίδιο αντικείμενο σε διαφορετικές γωνίες, καθώς τα διανύσματα χαρακτηριστικών θα είναι τελείως διαφορετικά. Ετσι δικαιολογείται η σημαντική πτώση της ακρίβειας.

Τέλος ο συνδυασμός των παραπάνω παραμορφώσεων υπογραμμίζει την διαφορά στην ανθεκτικότητα των δύο μεθόδων. Από την μία ο περιγραφητής \foreignlanguage{english}{SURF} διατηρεί παραπλήσια ποσοστά επιτυχίας με πριν ενώ από την άλλη ο \foreignlanguage{english}{HOG} διολισθαίνει σημαντικά φτάνοντας σε ποσοστό επιτυχίας 48\% από 61\%, δηλαδή για κάθε σωστή πρόβλεψη υπάρχει και μία λάθος. 

Αξίζει να σημειωθεί ότι στην σύγκριση ανάμεσα στους ανιχνευτές ο Hessian-\foreignlanguage{english}{Laplace} εμφανίζει με μικρή διαφορά αλλά καλύτερα ποσοστά έναντι του \foreignlanguage{english}{Harris-Laplace} σχεδόν σε όλες τις παραμορφώσεις. Αυτό θα μπορούσε να εξηγηθεί από το γεγονός ότι τα \foreignlanguage{english}{blobs} σαν σημεία ενδιαφέροντος αποτυπώνουν συμμετρικές δομές (άρα και αναλλοίωτες στην περιστροφή).

\begin{figure}[H] % Το [h] σημαίνει "here" (εδώ ακριβώς)
    %\centering    % Στοίχιση στο κέντρο
    \includegraphics[scale=0.5]{figs_part3/fig32.png}
    \caption{\foreignlanguage{english}{Accuracies and Standard Deviations Distorted Datasets - HOG SURF}} % Λεζάντα
    \label{fig:results} % Ετικέτα για αναφορά μέσα στο κείμενο (π.χ. Βλέπε Εικόνα \ref{fig:results})
\end{figure}


\section*{Παράρτημα}
\subsection*{Πίνακες Αποτελεσμάτων}
\hspace{0.6cm}Πειραματισμός με τις τιμές των παραμέτρων σ και θ, με τις μεθόδους \foreignlanguage{english}{Linear} και \foreignlanguage{english}{Non-Linear} για διαφορετική είσοδο (διαφορετικό \foreignlanguage{english}{PSNR} της εικόνας \foreignlanguage{english}{edgetest.jpg}). Παρουσιάζονται οι μετρικές \foreignlanguage{english}{Precision, Recall} και \foreignlanguage{english}{C-score}. Ο πίνακας παρατίθεται στο αρχείο \foreignlanguage{english}{133\_Results\_Matrix}. 




\subsection*{Εναλλακτική Προσέγγιση Παραγώγου Εξομαλυμένης Εικόνας}
\hspace{0.6cm}Εκμεταλλευόμενοι την ιδιότητα της συνέλιξης και των συστημάτων θα μπορούσαμε αντί για την χρήση της \foreignlanguage{english}{np.gradient} (που ακολουθεί μια διακριτή προσέγγιση της παραγώγου) να υπολογίσουμε τις αντίστοιχες αναλυτικές εκφράσεις των παραγώγων 1ης και 2ης τάξης ως προς \foreignlanguage{english}{x} και \foreignlanguage{english}{y} των Γκαουσιανών Φίλτρων και ύστερα να συνελίσσουμε με αυτά την αρχική εικόνα. Αυτή θα μας εξασφάλιζε μεγαλύτερη ακρίβεια, καλύτερα αποτελέσματα και μικρότερα σφάλματα προσέγγισης.

\section*{Πηγές}
Μαραγκός Π. Ανάλυση Εικόνων και Όραση Υπολογιστών. Ε.Μ.Π. 2014.


\end{document}

